\documentclass[10pt, letterpaper]{article}

% Inhaltsverzeichnis für Pakettypen (nur für Übersicht im Header, wird nicht im Dokument angezeigt)
% 1. Seitenlayout und Ränder
% 2. Sprache und Zeichensatz
% 3. Mathematik und Theorem-Umgebungen
% 4. Eigene Makros
% 5. Diagramme und Grafiken
% 6. Tabellen und Aufzählungen
% 7. Inhaltsverzeichnis
% 8. Abschnittsüberschriften
% 9. Abstrakt-Umgebung
% 10. Todos/Notizen
% 11. Rahmen/Box-Umgebungen
% 12. Python-Integration
% 13. Literaturverwaltung
% 14. Hyperlinks
% 15. Absatzeinstellungen
% 16. Umgebungen
% 17  Graphik
% 18  Extra
% 00. Titel und Autor

\usepackage{slashed}

% --- 1. Seitenlayout und Ränder ---
\usepackage[margin=3cm]{geometry}

% --- 2. Sprache und Zeichensatz ---
\usepackage[english]{babel}
\usepackage[T1]{fontenc}
\usepackage[utf8]{inputenc}

% --- 3. Mathematik und Theorem-Umgebungen ---
\usepackage{amsmath, amssymb, amsthm}
\usepackage{mathrsfs}
\DeclareMathOperator{\WF}{WF}

% --- 4. Eigene Makros ---
\usepackage{xcolor}
\newcommand{\SKP}{\langle\cdot,\cdot\rangle}
\newcommand{\id}{\operatorname{id}}
\newcommand{\sgn}{\operatorname{sgn}}
\newcommand{\Ad}{\operatorname{Ad}}
\newcommand{\ad}{\operatorname{ad}}
\newcommand{\R}{\mathbb{R}}
\newcommand{\N}{\mathbb{N}}
\newcommand{\Q}{\mathbb{Q}}
\newcommand{\Z}{\mathbb{Z}}
\newcommand{\C}{\mathbb{C}}
\newcommand{\QU}{\mathbb{H}}
\newcommand{\Cinf}{C^{\infty}}
\newcommand{\Mod}{\mathrm{Mod}}
\newcommand{\Alg}{\mathrm{Alg}}
\newcommand{\Diff}{\mathrm{Diff}}
\newcommand{\Iso}{\mathrm{Iso}}
\newcommand{\Aut}{\mathrm{Aut}}
\newcommand{\End}{\mathrm{End}}
\newcommand{\Hom}{\mathrm{Hom}}
\newcommand{\Cl}{\mathrm{Cl}}
\newcommand{\entwurf}[1]{\textcolor{red}{#1}}
\newcommand{\GL}{\mathrm{GL}}
\newcommand{\Mat}{\mathrm{Mat}}
\newcommand{\SL}{\mathrm{SL}}
\newcommand{\SO}{\mathrm{SO}}
\newcommand{\SU}{\mathrm{SU}}
\newcommand{\Sp}{\mathrm{Sp}}
\newcommand{\Spin}{\mathrm{Spin}}
\newcommand{\Pin}{\mathrm{Pin}}
\newcommand{\Lor}{\mathrm{Lor}}
\newcommand{\Ogroup}{\mathrm{O}}
\newcommand{\U}{\mathrm{U}}

% --- 5. Diagramme und Grafiken ---
\usepackage{graphicx}
\graphicspath{{../../Library/Images/}}
\usepackage[export]{adjustbox}
\usepackage{tikz}
\usetikzlibrary{decorations.pathreplacing, arrows.meta, positioning}
\usepackage{tikz-cd}

% --- 6. Tabellen und Aufzählungen ---
\usepackage{enumitem}
\setlist[itemize]{left=0.5cm}

\newenvironment{romanenum}[1][]
  {%
    \ifx&#1&
    \else
      \textbf{#1}\quad
    \fi
    \begin{enumerate}[label=\roman*)]
  }
  {%
    \end{enumerate}%
  }

% --- 7. Inhaltsverzeichnis ---
\usepackage{tocloft}
\renewcommand{\cftsecfont}{\footnotesize}
\renewcommand{\cftsubsecfont}{\footnotesize}
\renewcommand{\cftsubsubsecfont}{\footnotesize}
\renewcommand{\cftsecpagefont}{\footnotesize}
\renewcommand{\cftsubsecpagefont}{\footnotesize}
\renewcommand{\cftsubsubsecpagefont}{\footnotesize}
\usepackage{etoc}

% --- 8. Abschnittsüberschriften ---
\usepackage{titlesec}
\titleformat{\section}{\normalfont\large\bfseries}{\thesection}{1em}{}
\titleformat{\subsection}{\normalfont\normalsize\bfseries}{\thesubsection}{0.5em}{}
\titleformat{\subsubsection}{\normalfont\normalsize\bfseries}{\thesubsubsection}{0.5em}{}
\setcounter{secnumdepth}{4}

% --- 9. Abstrakt-Umgebung ---
\usepackage{changepage}
\renewenvironment{abstract}
  {
    \begin{adjustwidth}{1.5cm}{1.5cm}
    \small
    \textsc{Abstract. –}%
  }
  {
    \end{adjustwidth}
  }

% --- 10. Todos/Notizen ---
\usepackage{todonotes}
% Hellblau definieren
\definecolor{hellblau}{rgb}{0.3,0.6,1}
\newenvironment{note}{\par\begingroup\color{hellblau}\itshape}{\par\endgroup}

% --- 11. Rahmen/Box-Umgebungen ---
\usepackage{mdframed}
\usepackage{tcolorbox}
\colorlet{shadecolor}{gray!25}

\newenvironment{customTheorem}
  {\vspace{10pt}%
   \begin{mdframed}[
     backgroundcolor=gray!20,
     linewidth=0pt,
     innertopmargin=10pt,
     innerbottommargin=10pt,
     skipabove=\dimexpr\topsep+\ht\strutbox\relax,
     skipbelow=\topsep,
   ]}
  {\end{mdframed}
   \vspace{10pt}%
  }

% --- 12. Python-Integration ---
% (Deaktiviert in dieser Version, aktiviere bei Bedarf)
% \usepackage{pythontex}
% \usepackage[makestderr]{pythontex}

% --- 13. Literaturverwaltung ---
\usepackage{csquotes}
\usepackage[backend=biber, style=alphabetic, citestyle=alphabetic]{biblatex}
\addbibresource{bibliography.bib}

% --- 14. Hyperlinks ---
\usepackage{hyperref}
\hypersetup{
  colorlinks   = true,
  urlcolor     = blue,
  linkcolor    = blue,
  citecolor    = blue,
  frenchlinks  = true
}

% --- 15. Absatzeinstellungen ---
\usepackage[parfill]{parskip}
\sloppy

% --- 16. Umgebungen ---
\usepackage{thmtools}

\newcommand{\CustomHeading}[3]{%
  \par\medskip\noindent%
  \textbf{#1 #2} \textnormal{(#3)}.\enskip%
}

\newenvironment{DEF}[2]{\begin{unitbox}\CustomHeading{Definition}{#1}{#2}}{\end{unitbox}}
\newenvironment{PROP}[2]{\begin{unitbox}\CustomHeading{Proposition}{#1}{#2}}{\end{unitbox}}
\newenvironment{THEO}[2]{\begin{unitbox}\CustomHeading{Theorem}{#1}{#2}}{\end{unitbox}}
\newenvironment{LEM}[2]{\begin{unitbox}\CustomHeading{Lemma}{#1}{#2}}{\end{unitbox}}
\newenvironment{KORO}[2]{\begin{unitbox}\CustomHeading{Corollar}{#1}{#2}}{\end{unitbox}}
\newenvironment{REM}[2]{\begin{unitbox}\CustomHeading{Remark}{#1}{#2}}{\end{unitbox}}
\newenvironment{EXA}[2]{\begin{unitbox}\CustomHeading{Example}{#1}{#2}}{\end{unitbox}}
\newenvironment{STUD}[2]{\begin{unitbox}\CustomHeading{Study}{#1}{#2}}{\end{unitbox}}
\newenvironment{CONC}[2]{\begin{unitbox}\CustomHeading{Concept}{#1}{#2}}{\end{unitbox}}
\newenvironment{OTH}[2]{\begin{unitbox}\CustomHeading{Other}{#1}{#2}}{\end{unitbox}}
\newenvironment{EXE}[2]{\begin{unitbox}\CustomHeading{Exercise}{#1}{#2}}{\end{unitbox}}
\newenvironment{MOT}[2]{\begin{unitbox}\CustomHeading{Motivation}{#1}{#2}}{\end{unitbox}}
\newenvironment{PROOF}[2]{\begin{unitbox}\CustomHeading{Proof}{#1}{#2}}{\end{unitbox}}

% --- Unit Umgebung für Source-Inhalte ---
\usepackage{mdframed}
\newmdenv[
  linewidth=1pt,
  topline=false,
  bottomline=false,
  rightline=false,
  leftmargin=0cm,
  rightmargin=0cm,
  skipabove=10pt,
  skipbelow=10pt,
  innertopmargin=0.5\baselineskip,
  innerbottommargin=0.5\baselineskip,
  backgroundcolor=gray!10,
  linecolor=gray
]{unitbox}

\newenvironment{unit}[1]
  {\begin{unitbox}\textbf{Unit #1}\par\smallskip}
  {\end{unitbox}}

% --- 17. ... ---

% --- 18. Extras ---
\usepackage{stmaryrd}
\usepackage{bbold}  % falls du \mathbb{1} nutzen willst

% --- 00. Titel und Autor ---
\title{QFT Survey}
\author{Tim Jaschik}
\date{\today}

\begin{document}

\maketitle
\rule{\textwidth}{0.5pt}
\begin{abstract}
Kurze Beschreibung …
\end{abstract}
\rule{\textwidth}{0.5pt}
\vspace{0.5cm}

\tableofcontents

\pagebreak

\section{Notes}


\begin{enumerate}
  \item Kinematik: Heisenbergalgebra, Observablen und Darstellungen
  \begin{enumerate}
    \item Heisenberg Algebra
    \item Observablen
    \item Kommutatorrelationen
    \item Topologische Vektorräume von temperierten Distributionen und unitär äquivalente Darstellungen
    \item Fouriertransformation
    \item 
  \end{enumerate} 
  \item Dynamik: 
  \begin{enumerate}
    \item Heisenberg Bild (Dynamik wirkt auf Observablen)
    \begin{enumerate}
      \item Hüllalgebra der Heisenbergalgebra  
    \end{enumerate}
    \item Schrödinger Bild (Dynamik wirkt auf Zustände)
    \begin{enumerate}
      \item 
    \end{enumerate}
    \item Auswahlverfahren / Kriterien für interessante Hamiltonians
  \end{enumerate}
  \item Symmetrien in Theoriebildung
  \begin{enumerate}
  \item Mathematische Invarainten-Forderungen und Lie Formalismus
  \item Lie Theorie in QM
  \item SRT und Kovarainzbedingungen an Theorien
  \begin{enumerate}
    \item SRT
    \item Lorentzgruppe und Poincare Gruppe
    \item Konzept der Kovarianz von Theorien
    \item Interpretation von Lorentztransformationen
  \end{enumerate}
  \end{enumerate}
  \item Quantisierung von klassischer Feldtheorie
  \begin{enumerate}
    \item Konstruktion von Fock-Raum
  \end{enumerate}
\end{enumerate}

\pagebreak


\section{Quantenmechanische Theorien}

\begin{enumerate}
  \item Erwartungshaltung an quantenmechanische Theorien
  \item Axiomatiken
  \item Allgemeiner vs. spezifischer Rahmen (QFT)
\end{enumerate}


\begin{unit}
\emph{Zwischenebene: Abstrakte Heisenberg--/Weyl-Kinematik für QM und QFT}

In diesem Abschnitt formulieren wir einen abstrakten Rahmen, der
\emph{sowohl} die klassische Quantenmechanik mit endlich vielen
Freiheitsgraden \((X=\mathbb{R}^{2n})\),
\emph{als auch} freie (bosonische oder fermionische) Quantenfeldtheorien
\((X\ \text{unendlichdimensional})\)
und sogar die üblichen einfachen Wechselwirkungen der QFT1/2
gemeinsam erfasst.
Diese Ebene liegt strenger und allgemeiner als die kanonische
Heisenberg-/Weyl-Algebra über \(\mathbb{R}^{2n}\),
aber konkreter als der vollabstrakte C$^\ast$-/von-Neumann-Formalismus.
Wir nennen sie die \emph{Zwischenebene der Einteilchenräume und CCR/CAR-Algebren}.

%---------------------------------------------------------
\emph{1. Einteilchenraum und symplektische/skalare Struktur}

\emph{Bosonische Systeme.}
Sei \(X\) ein reeller lokalkonvexer Vektorraum
(oft endlichdimensional, oder ein Testfunktionenraum wie \(\mathcal{S}(\mathbb{R}^d)\)),
versehen mit einer nichtausgearteten, antisymmetrischen Bilinearform
\[
\sigma : X\times X \to \mathbb{R}.
\]
Das Paar \((X,\sigma)\) heißt \emph{symplektischer Einteilchenraum}.
\begin{itemize}
  \item QM: \(X=\mathbb{R}^{2n}\) mit der kanonischen Form \(\sigma((x,p),(x',p'))=p\cdot x'-p'\cdot x\).
  \item Freie QFT: \(X\) ist der Raum reeller Testfunktionen oder Cauchy-Daten;
        \(\sigma\) ist z.\,B. durch das Pauli–Jordan-Funktional gegeben.
\end{itemize}

\emph{Fermionische Systeme.}
Sei \(Y\) ein reeller Hilbertraum
mit innerem Produkt \(\langle\cdot,\cdot\rangle_Y\).
Dies ist der \emph{Einteilchenraum für Fermionen}.
\begin{itemize}
  \item QM (Spin): \(Y=\mathbb{R}^n\), inneres Produkt gewöhnlich.
  \item Freies Dirac- oder Majorana-Feld: \(Y\) ist ein Raum glatter kompakter Spinorfunktionen.
\end{itemize}

Diese Wahl \((X,\sigma)\) bzw. \((Y,\langle\cdot,\cdot\rangle)\) ist der
\emph{kinematische Kern} der Theorie: sie entspricht dem klassischen
Phasenraum (bosonisch) bzw. klassischen Konfigurationsraum (fermionisch).

%---------------------------------------------------------
\emph{2. Abstrakte Heisenberg-Lie-Algebra}

Zu einem symplektischen Einteilchenraum \((X,\sigma)\) definieren wir:

\begin{DEF}{MID-HEIS}{Abstrakte Heisenberg-Algebra}
\[
\mathfrak{h}(X) := X \oplus \mathbb{R}Z
\]
mit Lie-Klammer
\[
[(f,\alpha Z),(g,\beta Z)]
= (0,\ \sigma(f,g)\,Z),
\qquad [Z,\cdot]=0.
\]
\end{DEF}

Ein Element \(f\in X\) wird physikalisch als „abstrakter Feldmodus“
oder „abstrakte Kanonische Observable“ interpretiert;
\(Z\) spielt die Rolle von \(i\hbar\).

Dies umfasst:
\begin{itemize}
  \item QM: \(\mathfrak{h}(\mathbb{R}^{2n})\) ist die übliche Heisenberg-Algebra
        mit endlich vielen \(X_j,P_k\).
  \item Freie QFT: \(\mathfrak{h}(X)\) ist unendlichdimensional und kodiert
        \([\,\phi(f),\pi(g)\,]= i\int f g\).
\end{itemize}

Zu \(\mathfrak{h}(X)\) gehört die universelle Hüllalgebra
\(U(\mathfrak{h}(X))\), die die polynomiellen Kombinationen dieser
abstrakten Observablen enthält.

%---------------------------------------------------------
\emph{3. CCR- und CAR-C$^\ast$-Algebren (Weyl- und Clifford-Ebene)}

Die Heisenberg-Algebra ist nur die infinitesimale Struktur.
Die physikalisch sinnvolle Observablenalgebra für QM und freie QFT
ist die exponentialisierte C$^\ast$-Algebra.

\begin{DEF}{MID-CCR}{CCR-Weyl-C$^\ast$-Algebra (bosonisch)}
Für \((X,\sigma)\) ist die \emph{Weyl- bzw. CCR-C$^\ast$-Algebra}
\(\mathrm{CCR}(X,\sigma)\)
die vollständige C$^\ast$-Algebra, erzeugt durch Symbole \(W(f)\) (\(f\in X\))
mit den Weyl-Relationen
\[
W(f)W(g)=e^{-\,\frac{i}{2}\sigma(f,g)}\,W(f+g),\qquad W(f)^\ast=W(-f).
\]
\end{DEF}

\begin{DEF}{MID-CAR}{CAR-C$^\ast$-Algebra (fermionisch)}
Für einen reellen Hilbertraum \(Y\) ist die \emph{CAR-Algebra}
\(\mathrm{CAR}(Y)\) die C$^\ast$-Algebra, erzeugt durch Symbole \(a(f)\) (\(f\in Y\))
mit den Relationen
\[
\{a(f),a(g)\}=0,\qquad
\{a(f),a(g)^\ast\}=\langle f,g\rangle_Y\mathbf{1}.
\]
\end{DEF}

Diese Konstruktionen umfassen:
\begin{itemize}
  \item gewöhnliche QM (\(X\) endlichdimensional),
  \item freie skalare und elektromagnetische Felder (CCR),
  \item freie Dirac- und Majorana-Felder (CAR),
  \item alle Standardbeispiele aus QFT1/2 auf einer gemeinsamen algebraischen Ebene.
\end{itemize}

%---------------------------------------------------------
\emph{4. Dynamik als Automorphismengruppe der CCR/CAR-Algebra}

\begin{DEF}{MID-DYN}{Abstrakte Dynamik}
Eine (Heisenberg-)Dynamik ist eine stark-stetige Gruppe von *-Automorphismen
\[
\alpha : \mathbb{R} \to \mathrm{Aut}(\mathcal{A}),\qquad
t\mapsto \alpha_t,
\]
wobei \(\mathcal{A}=\mathrm{CCR}(X,\sigma)\) oder \(\mathrm{CAR}(Y)\).
\end{DEF}

\begin{itemize}
  \item \textbf{Freie Theorien:}  
        Die Einteilchen-Dynamik \(T_t:X\to X\) (oder \(Y\to Y\))
        induziert eine Automorphismengruppe
        \[
        \alpha_t(W(f)) = W(T_tf),\qquad
        \alpha_t(a(f)) = a(T_t f).
        \]
  \item \textbf{Wechselwirkungen in QFT1/2 (formal):}  
        dieselbe Algebra bleibt gültig, aber die Zeitentwicklung wird
        durch einen formalen Hamiltonoperator \(H=H_0+V\) bestimmt:
        \[
        \frac{d}{dt}\alpha_t(A)=i[H,A].
        \]
        Die Algebra ändert sich also nicht – nur die Dynamik.
\end{itemize}

Damit lassen sich alle Standardmodelle einer QFT-Einführung
(Klein–Gordon, Dirac, Yukawa, \(\phi^4\), QED in der kanonischen Form)
im selben algebraischen Rahmen behandeln.

%---------------------------------------------------------
\emph{5. Darstellungen (GNS) und Fockräume}

Ein Zustand \(\omega\) auf \(\mathrm{CCR}(X,\sigma)\) oder \(\mathrm{CAR}(Y)\)
definiert über die GNS-Konstruktion eine Darstellung
\[
(\pi_\omega,\mathcal{H}_\omega,\Omega_\omega).
\]

\begin{itemize}
  \item Für freie Theorien sind \(\omega\) \emph{quasi-freie Zustände},
        und die GNS-Darstellung ist die übliche \emph{bosonische} bzw.
        \emph{fermionische Fock-Darstellung}.
  \item Für Wechselwirkungen in QFT1/2 bleibt formal dieselbe Algebra bestehen,
        aber der Zustand \(\omega_V\) (interacting vacuum) ist i.\,Allg.
        nichtunitär äquivalent zum freien Zustand \(\omega_0\).
\end{itemize}

%---------------------------------------------------------
\emph{6. Charakterisierung dieser Zwischenebene}

\begin{REM}{MID-SUM}{Zusammenfassung}
Die hier eingeführte Ebene
\[
(X,\sigma)\longmapsto \mathrm{CCR}(X,\sigma),\qquad
Y\longmapsto\mathrm{CAR}(Y)
\]
liefert einen abstrakten Rahmen, der
\begin{itemize}
  \item die kinematische Struktur der QM 
        (endlichdimensionaler symplektischer Raum),
  \item die kinematische Struktur aller freien Feldtheorien
        (unendlichdimensionale Einteilchenräume),
  \item sowie die formale Behandlung einfacher/interagierender QFT-Modelle
        (gleiche Algebra, veränderte Dynamik/Zustände)
\end{itemize}
einheitlich beschreibt.

Diese Ebene ist damit:
\begin{enumerate}
  \item \emph{allgemeiner} als die klassische Heisenberg-/Weyl-Algebra über \(\mathbb{R}^{2n}\),
  \item \emph{konkret genug}, um die vollständige Standard-QFT1/2-Abdeckung zu gewährleisten,
  \item \emph{kompatibel} mit dem noch allgemeineren C$^\ast$-/von-Neumann-Formalismus,
        in dem CCR/CAR-Algebren als strukturierte Spezialfälle erscheinen.
\end{enumerate}
\end{REM}
\end{unit}


\begin{unit}
\emph{Endliche Quantenmechanik als Spezialfall der abstrakten CCR-Struktur}

In diesem Abschnitt zeigen wir explizit, wie die gewöhnliche nichtrelativistische
Quantenmechanik mit endlich vielen Freiheitsgraden in den abstrakten Rahmen
\[
(X,\sigma)\;\longmapsto\;\mathrm{CCR}(X,\sigma)
\]
eingebettet ist. Ziel ist es, die üblichen Objekte
\[
\mathcal{H}=L^2(\mathbb{R}^n),\quad
\hat{x}_j,\hat{p}_k,\quad
H=\tfrac{1}{2m}\hat{p}^2+V(\hat{x})
\]
als konkrete Realisierung der abstrakten Heisenberg- und CCR-Struktur zu identifizieren.

%--------------------------------------------------------------------
\emph{1. Kinematik: Der symplektische Einteilchenraum}

\emph{Abstrakt.}
Sei \(V\simeq\mathbb{R}^n\) der klassische Konfigurationsraum.
Der \emph{klassische Phasenraum} ist \(T^\ast V\simeq \mathbb{R}^{2n}\).
Wir setzen
\[
X := V_x \oplus V_p \;\cong\; \mathbb{R}^n_x\oplus\mathbb{R}^n_p\cong\mathbb{R}^{2n}
\]
mit der kanonischen symplektischen Form
\[
\sigma\bigl((x,p),(x',p')\bigr)
:= p\cdot x' - p'\cdot x.
\]
Damit ist \((X,\sigma)\) ein endlichdimensionaler symplektischer Raum.

\emph{Interpretation.}
Die Paare \((x,p)\in X\) repräsentieren klassische Zustände
\((\text{Ort},\text{Impuls})\), und \(\sigma\) ist die übliche symplektische Struktur
des Phasenraums.

%--------------------------------------------------------------------
\emph{2. Abstrakte Heisenberg-Algebra \texorpdfstring{$\mathfrak{h}(X)$}{h(X)}}

Nach der allgemeinen Definition (vgl.\ \textsc{MID-HEIS}) ist die
\emph{Heisenberg-Lie-Algebra}
\[
\mathfrak{h}(X) := X \oplus \mathbb{R}Z
\]
mit Lie-Klammer
\[
[(f,\alpha Z),(g,\beta Z)]
= (0,\ \sigma(f,g)\,Z).
\]
Wählt man eine Basis \(\{e_j\}_{j=1}^n\) von \(V_x\) und \(\{e^j\}_{j=1}^n\) von \(V_p\),
so kann man
\[
X_j := (e_j,0),\qquad
P_j := (0,e^j)
\]
als Erzeuger wählen. Dann lauten die Relationen
\[
[X_j,P_k] = \delta_{jk}\,Z,\qquad
[X_j,X_k]=[P_j,P_k]=[Z,\cdot]=0.
\]
Identifiziert man formal \(Z\sim i\hbar\), erhält man die bekannten
kanonischen Kommutationsrelationen:
\[
[X_j,P_k] = i\hbar\,\delta_{jk}.
\]

Die universelle Hüllalgebra \(U(\mathfrak{h}(X))\) enthält dann
Polynome in den \(X_j,P_k\), z.\,B. Hamiltonoperatoren der Form
\[
H = \frac{1}{2m}\sum_{j=1}^n P_j^2 + V(X_1,\dots,X_n).
\]

%--------------------------------------------------------------------
\emph{3. Die CCR-Weyl-C$^\ast$-Algebra \texorpdfstring{$\mathrm{CCR}(X,\sigma)$}{CCR(X,sigma)}}

\emph{Abstrakte Definition.}
Die CCR- (oder Weyl-) C$^\ast$-Algebra \(\mathrm{CCR}(X,\sigma)\)
ist die C$^\ast$-Vervollständigung der *-Algebra,
erzeugt von Symbolen \(W(f)\) (\(f\in X\)) mit den Weyl-Relationen
\[
W(f)W(g) = e^{-\frac{i}{2}\sigma(f,g)}W(f+g),\qquad
W(f)^\ast = W(-f).
\]

Diese Algebra ist ein abstrakter Vertreter der kanonischen
Kommutationsrelationen. Die \emph{infinitesimale} Heisenberg-Algebra
\(\mathfrak{h}(X)\) kann als Lie-Algebra der „logarithmischen Ableitungen“
der \(W(f)\) verstanden werden.

%--------------------------------------------------------------------
\emph{4. Schrödinger-Darstellung als konkrete Realisierung}

\emph{Hilbertraum.}
Wir wählen als Zustandsraum den Hilbertraum
\[
\mathcal{H} := L^2(\mathbb{R}^n,dx),
\]
mit kanonischem Skalarprodukt
\[
\langle \psi,\varphi\rangle = \int_{\mathbb{R}^n} \overline{\psi(x)}\varphi(x)\,dx.
\]

\emph{Ort- und Impulsoperatoren.}
Wir definieren (zunächst auf einem dichten Unterraum, etwa den Schwartzfunktionen)
\[
(\hat{X}_j\psi)(x) := x_j\psi(x),\qquad
(\hat{P}_j\psi)(x) := -\,i\hbar\,\frac{\partial}{\partial x_j}\psi(x).
\]
Dann gilt (formal)
\[
[\hat{X}_j,\hat{P}_k]\psi
= i\hbar\,\delta_{jk}\,\psi,\qquad
[\hat{X}_j,\hat{X}_k]=[\hat{P}_j,\hat{P}_k]=0,
\]
d.\,h. \(\hat{X}_j,\hat{P}_k\) realisieren die Relationen der Heisenberg-Algebra
\(\mathfrak{h}(X)\).

\emph{Weyl-Operatoren in der Schrödinger-Darstellung.}
Zu jedem \(f=(x,p)\in X\) definieren wir einen unitären Operator
\(\pi_{\mathrm{Sch}}(W(f))\) auf \(\mathcal{H}\) durch
\[
\bigl(\pi_{\mathrm{Sch}}(W(x,p))\psi\bigr)(y)
:= e^{\frac{i}{\hbar}\left(p\cdot y - \tfrac{1}{2}p\cdot x\right)}\,
   \psi(y-x).
\]

Man überprüft:
\begin{align*}
\pi_{\mathrm{Sch}}(W(x,p))\pi_{\mathrm{Sch}}(W(x',p'))
&= e^{-\frac{i}{2\hbar}\sigma((x,p),(x',p'))}\,
   \pi_{\mathrm{Sch}}(W(x+x',p+p')),\\
\pi_{\mathrm{Sch}}(W(x,p))^\ast
&= \pi_{\mathrm{Sch}}(W(-x,-p)),
\end{align*}
also erfüllt \(\pi_{\mathrm{Sch}}\) genau die Weyl-Relationen.
Damit ist
\[
\pi_{\mathrm{Sch}}:\mathrm{CCR}(X,\sigma)\longrightarrow \mathcal{B}(\mathcal{H})
\]
eine \emph{konkrete Darstellung} der abstrakten CCR-Algebra.

\emph{Infinitesimale Generatoren.}
Differenziert man \(t\mapsto \pi_{\mathrm{Sch}}(W(t e_j,0))\) und
\(t\mapsto \pi_{\mathrm{Sch}}(W(0,t e^j))\) bei \(t=0\), so erhält man
\[
\hat{X}_j = i\hbar\,\frac{d}{dt}\pi_{\mathrm{Sch}}(W(t e_j,0))\bigg|_{t=0},
\quad
\hat{P}_j = i\hbar\,\frac{d}{dt}\pi_{\mathrm{Sch}}(W(0,t e^j))\bigg|_{t=0}.
\]
Damit erscheinen die üblichen Ort- und Impulsoperatoren
als \emph{infinitesimale Generatoren} der Weyl-Darstellung
von \(\mathrm{CCR}(X,\sigma)\).

%--------------------------------------------------------------------
\emph{5. Stone--von Neumann und Eindeutigkeit (bis auf Äquivalenz)}

\begin{THEO}{SvN-1}{Stone--von Neumann (informell)}
Sei \(\pi:\mathrm{CCR}(X,\sigma)\to \mathcal{B}(\mathcal{H}_\pi)\)
eine irreduzible, \emph{regular} (d.\,h. stetige) Darstellung
der CCR-Algebra für endlichdimensionales \(X\).
Dann ist \(\pi\) zu der Schrödinger-Darstellung \(\pi_{\mathrm{Sch}}\)
unitär äquivalent.
\end{THEO}

\emph{Interpretation.}
Für eine gegebene endliche Anzahl von Freiheitsgraden \(n\) ist die
Struktur der Hilbertraumdarstellung der Heisenberg-/Weyl-Algebra
durch \(L^2(\mathbb{R}^n)\) (bis auf unitäre Äquivalenz) \emph{eindeutig}.
Damit ist der abstrakte Rahmen \((X,\sigma)\mapsto\mathrm{CCR}(X,\sigma)\)
für die endliche QM nicht nur realisierbar, sondern im Wesentlichen
\emph{vollständig}.

%--------------------------------------------------------------------
\emph{6. Hamiltonoperatoren und Dynamik in \texorpdfstring{$U(\mathfrak{h}(X))$}{U(h(X))}}

\emph{Hamiltonoperatoren als Elemente der Hüllalgebra.}
In der abstrakten Hüllalgebra \(U(\mathfrak{h}(X))\) können wir
Hamiltonoperatoren als Polynome in \(X_j,P_k\) definieren, z.\,B.
\[
H = \frac{1}{2m}\sum_{j=1}^n P_j^2 + V(X_1,\dots,X_n).
\]
In der Schrödinger-Darstellung wird daraus der übliche Operator
\[
\hat{H} = -\frac{\hbar^2}{2m}\Delta + V(x)
\]
auf \(L^2(\mathbb{R}^n)\).

\emph{Dynamik.}
Die Heisenberg-Zeitentwicklung eines Observablenoperators \(A\) wird
abstrakt durch die innere Derivation
\[
\dot{A}(t) = \frac{i}{\hbar}[H,A(t)]
\]
in \(U(\mathfrak{h}(X))\) beschrieben.
In der Schrödinger-Darstellung entspricht dies der üblichen Formel
\[
A_H(t) = e^{\frac{i}{\hbar}t\hat{H}} A_S e^{-\frac{i}{\hbar}t\hat{H}},
\]
während die Zustandsdynamik (Schrödinger-Bild) durch
\(i\hbar\partial_t\psi = \hat{H}\psi\) gegeben ist.

Damit ist die volle Dynamik der nichtrelativistischen QM
ein Spezialfall der abstrakten Dynamik in \(U(\mathfrak{h}(X))\)
und ihrer Darstellung auf \(L^2(\mathbb{R}^n)\).

%--------------------------------------------------------------------
\emph{7. Fazit: Endliche QM als Spezialfall der Zwischenebene}

\begin{REM}{QM-SPECIAL}{Zusammenfassung}
Die gewöhnliche Quantenmechanik mit endlich vielen Freiheitsgraden
entsteht aus der Zwischenebene
\[
(X,\sigma)\longmapsto \mathrm{CCR}(X,\sigma)
\]
spezifisch durch die Wahl:
\begin{itemize}
  \item \(X = \mathbb{R}^{2n}\) mit kanonischer symplektischer Form \(\sigma\).
  \item CCR-Weyl-C$^\ast$-Algebra \(\mathcal{A} = \mathrm{CCR}(X,\sigma)\) als Observablenalgebra.
  \item Schrödinger-Darstellung \(\pi_{\mathrm{Sch}}:\mathcal{A}\to\mathcal{B}(L^2(\mathbb{R}^n))\),
        die nach Stone--von Neumann (unter Regularität) bis auf unitäre Äquivalenz eindeutig ist.
  \item Heisenberg-Algebra \(\mathfrak{h}(X)\) und Hüllalgebra \(U(\mathfrak{h}(X))\),
        deren Generatoren sich als die üblichen Ort- und Impulsoperatoren
        realisieren.
  \item Hamiltonoperatoren \(H \in U(\mathfrak{h}(X))\) und die daraus
        abgeleitete Dynamik \(\dot{A}=\frac{i}{\hbar}[H,A]\) im Heisenberg-Bild
        bzw.\ \(i\hbar\partial_t\psi=\hat{H}\psi\) im Schrödinger-Bild.
\end{itemize}

In diesem präzisen Sinn ist die endliche nichtrelativistische QM
ein \emph{Spezialfall} der abstrakten Heisenberg-/CCR-Struktur:
die allgemeineren Konstruktionen (symplektischer Einteilchenraum,
Heisenberg-Algebra, CCR-Algebra, Automorphismengruppe)
reduzieren sich für \(X=\mathbb{R}^{2n}\) genau auf das übliche
Quantensystem \((L^2(\mathbb{R}^n),\hat{x},\hat{p},\hat{H})\).
\end{REM}


\end{unit}


\section{Kinematik einer Quantenmechanischen Theorie}


\begin{enumerate}
  \item Heisenberg Algebra
  \item Elementare Observablen
  \item Zustandsräume
\end{enumerate}


\section{Dynamik einer Quantenmechanischen Theorie}

\begin{enumerate}
  \item Heisenberg Bild
  \item Weyl Algebra und Polynome in elementaren Observablen
  \item Auswahlverfahren / Kriterien für interessante Hamiltonians
  \item Schrödinger Bild
\end{enumerate}


\section{Symmetrien in Theoriebildung}

  \begin{enumerate}
  \item Physikalisch motivierte Invarainten-Forderungen und Lie Formalismus
  \item Physikalische Lie Gruppen (Lorentz, Poincare, Galilei, Pin, Spin, SO, SU, U)
  \item Lie Theorie in QM
  \item SRT und Kovarainzbedingungen an Theorien
  \begin{enumerate}
    \item SRT
    \item Lorentzgruppe 
    \item Poincare Gruppe
    \item Konzept der Kovarianz von Theorien
    \item Interpretation von Lorentztransformationen
  \end{enumerate}
  \item CPT Trafos
  \end{enumerate}


\section{Relativistische Wellengleichungsmodelle}

\begin{enumerate}
  \item Rolle von RWG-Modellen
  \item Beziehung zu quantisierten Feldtheorien
  \item Unitär äquivalente Darstellungen und Fourier Trafo
\end{enumerate}


\subsection{Klein Gordon}



\subsection{Dirac}

\begin{enumerate}
  \item Dirac Gleichungen (Schrödinger / Kovariant / Adjungiert / Charge Konjugiert)
  \item Parameter Charakterisierung
  \item Kovarianz bzgl $\Spin^+(1,3)$
  \item Lösungsentwicklung
  \item Verträglichkeit mit Grenzbereichen des Parameterraumes
  \item Nichtrelativistischer Bereich
  \item Ultrarelativistischer Bereich
\end{enumerate}



\section{Quantisierung von Feldtheorien}

\begin{enumerate}
  \item Motivation für Quantisierung von Feldtheorien
\end{enumerate}


\pagebreak

\section{Body}


\subsection{Abstraktes Setting}

\begin{unit}
\emph{Zwischenebene: Abstrakte Heisenberg--/Weyl-Kinematik für QM und QFT}

In diesem Abschnitt formulieren wir einen abstrakten Rahmen, der
\emph{sowohl} die klassische Quantenmechanik mit endlich vielen
Freiheitsgraden \((X=\mathbb{R}^{2n})\),
\emph{als auch} freie (bosonische oder fermionische) Quantenfeldtheorien
\((X\ \text{unendlichdimensional})\)
und sogar die üblichen einfachen Wechselwirkungen der QFT1/2
gemeinsam erfasst.
Diese Ebene liegt strenger und allgemeiner als die kanonische
Heisenberg-/Weyl-Algebra über \(\mathbb{R}^{2n}\),
aber konkreter als der vollabstrakte C$^\ast$-/von-Neumann-Formalismus.
Wir nennen sie die \emph{Zwischenebene der Einteilchenräume und CCR/CAR-Algebren}.

%---------------------------------------------------------
\emph{1. Einteilchenraum und symplektische/skalare Struktur}

\emph{Bosonische Systeme.}
Sei \(X\) ein reeller lokalkonvexer Vektorraum
(oft endlichdimensional, oder ein Testfunktionenraum wie \(\mathcal{S}(\mathbb{R}^d)\)),
versehen mit einer nichtausgearteten, antisymmetrischen Bilinearform
\[
\sigma : X\times X \to \mathbb{R}.
\]
Das Paar \((X,\sigma)\) heißt \emph{symplektischer Einteilchenraum}.
\begin{itemize}
  \item QM: \(X=\mathbb{R}^{2n}\) mit der kanonischen Form \(\sigma((x,p),(x',p'))=p\cdot x'-p'\cdot x\).
  \item Freie QFT: \(X\) ist der Raum reeller Testfunktionen oder Cauchy-Daten;
        \(\sigma\) ist z.\,B. durch das Pauli–Jordan-Funktional gegeben.
\end{itemize}

\emph{Fermionische Systeme.}
Sei \(Y\) ein reeller Hilbertraum
mit innerem Produkt \(\langle\cdot,\cdot\rangle_Y\).
Dies ist der \emph{Einteilchenraum für Fermionen}.
\begin{itemize}
  \item QM (Spin): \(Y=\mathbb{R}^n\), inneres Produkt gewöhnlich.
  \item Freies Dirac- oder Majorana-Feld: \(Y\) ist ein Raum glatter kompakter Spinorfunktionen.
\end{itemize}

Diese Wahl \((X,\sigma)\) bzw. \((Y,\langle\cdot,\cdot\rangle)\) ist der
\emph{kinematische Kern} der Theorie: sie entspricht dem klassischen
Phasenraum (bosonisch) bzw. klassischen Konfigurationsraum (fermionisch).

%---------------------------------------------------------
\emph{2. Abstrakte Heisenberg-Lie-Algebra}

Zu einem symplektischen Einteilchenraum \((X,\sigma)\) definieren wir:

\begin{DEF}{MID-HEIS}{Abstrakte Heisenberg-Algebra}
\[
\mathfrak{h}(X) := X \oplus \mathbb{R}Z
\]
mit Lie-Klammer
\[
[(f,\alpha Z),(g,\beta Z)]
= (0,\ \sigma(f,g)\,Z),
\qquad [Z,\cdot]=0.
\]
\end{DEF}

Ein Element \(f\in X\) wird physikalisch als „abstrakter Feldmodus“
oder „abstrakte Kanonische Observable“ interpretiert;
\(Z\) spielt die Rolle von \(i\hbar\).

Dies umfasst:
\begin{itemize}
  \item QM: \(\mathfrak{h}(\mathbb{R}^{2n})\) ist die übliche Heisenberg-Algebra
        mit endlich vielen \(X_j,P_k\).
  \item Freie QFT: \(\mathfrak{h}(X)\) ist unendlichdimensional und kodiert
        \([\,\phi(f),\pi(g)\,]= i\int f g\).
\end{itemize}

Zu \(\mathfrak{h}(X)\) gehört die universelle Hüllalgebra
\(U(\mathfrak{h}(X))\), die die polynomiellen Kombinationen dieser
abstrakten Observablen enthält.

%---------------------------------------------------------
\emph{3. CCR- und CAR-C$^\ast$-Algebren (Weyl- und Clifford-Ebene)}

Die Heisenberg-Algebra ist nur die infinitesimale Struktur.
Die physikalisch sinnvolle Observablenalgebra für QM und freie QFT
ist die exponentialisierte C$^\ast$-Algebra.

\begin{DEF}{MID-CCR}{CCR-Weyl-C$^\ast$-Algebra (bosonisch)}
Für \((X,\sigma)\) ist die \emph{Weyl- bzw. CCR-C$^\ast$-Algebra}
\(\mathrm{CCR}(X,\sigma)\)
die vollständige C$^\ast$-Algebra, erzeugt durch Symbole \(W(f)\) (\(f\in X\))
mit den Weyl-Relationen
\[
W(f)W(g)=e^{-\,\frac{i}{2}\sigma(f,g)}\,W(f+g),\qquad W(f)^\ast=W(-f).
\]
\end{DEF}

\begin{DEF}{MID-CAR}{CAR-C$^\ast$-Algebra (fermionisch)}
Für einen reellen Hilbertraum \(Y\) ist die \emph{CAR-Algebra}
\(\mathrm{CAR}(Y)\) die C$^\ast$-Algebra, erzeugt durch Symbole \(a(f)\) (\(f\in Y\))
mit den Relationen
\[
\{a(f),a(g)\}=0,\qquad
\{a(f),a(g)^\ast\}=\langle f,g\rangle_Y\mathbf{1}.
\]
\end{DEF}

Diese Konstruktionen umfassen:
\begin{itemize}
  \item gewöhnliche QM (\(X\) endlichdimensional),
  \item freie skalare und elektromagnetische Felder (CCR),
  \item freie Dirac- und Majorana-Felder (CAR),
  \item alle Standardbeispiele aus QFT1/2 auf einer gemeinsamen algebraischen Ebene.
\end{itemize}

%---------------------------------------------------------
\emph{4. Dynamik als Automorphismengruppe der CCR/CAR-Algebra}

\begin{DEF}{MID-DYN}{Abstrakte Dynamik}
Eine (Heisenberg-)Dynamik ist eine stark-stetige Gruppe von *-Automorphismen
\[
\alpha : \mathbb{R} \to \mathrm{Aut}(\mathcal{A}),\qquad
t\mapsto \alpha_t,
\]
wobei \(\mathcal{A}=\mathrm{CCR}(X,\sigma)\) oder \(\mathrm{CAR}(Y)\).
\end{DEF}

\begin{itemize}
  \item \textbf{Freie Theorien:}  
        Die Einteilchen-Dynamik \(T_t:X\to X\) (oder \(Y\to Y\))
        induziert eine Automorphismengruppe
        \[
        \alpha_t(W(f)) = W(T_tf),\qquad
        \alpha_t(a(f)) = a(T_t f).
        \]
  \item \textbf{Wechselwirkungen in QFT1/2 (formal):}  
        dieselbe Algebra bleibt gültig, aber die Zeitentwicklung wird
        durch einen formalen Hamiltonoperator \(H=H_0+V\) bestimmt:
        \[
        \frac{d}{dt}\alpha_t(A)=i[H,A].
        \]
        Die Algebra ändert sich also nicht – nur die Dynamik.
\end{itemize}

Damit lassen sich alle Standardmodelle einer QFT-Einführung
(Klein–Gordon, Dirac, Yukawa, \(\phi^4\), QED in der kanonischen Form)
im selben algebraischen Rahmen behandeln.

%---------------------------------------------------------
\emph{5. Darstellungen (GNS) und Fockräume}

Ein Zustand \(\omega\) auf \(\mathrm{CCR}(X,\sigma)\) oder \(\mathrm{CAR}(Y)\)
definiert über die GNS-Konstruktion eine Darstellung
\[
(\pi_\omega,\mathcal{H}_\omega,\Omega_\omega).
\]

\begin{itemize}
  \item Für freie Theorien sind \(\omega\) \emph{quasi-freie Zustände},
        und die GNS-Darstellung ist die übliche \emph{bosonische} bzw.
        \emph{fermionische Fock-Darstellung}.
  \item Für Wechselwirkungen in QFT1/2 bleibt formal dieselbe Algebra bestehen,
        aber der Zustand \(\omega_V\) (interacting vacuum) ist i.\,Allg.
        nichtunitär äquivalent zum freien Zustand \(\omega_0\).
\end{itemize}

%---------------------------------------------------------
\emph{6. Charakterisierung dieser Zwischenebene}

\begin{REM}{MID-SUM}{Zusammenfassung}
Die hier eingeführte Ebene
\[
(X,\sigma)\longmapsto \mathrm{CCR}(X,\sigma),\qquad
Y\longmapsto\mathrm{CAR}(Y)
\]
liefert einen abstrakten Rahmen, der
\begin{itemize}
  \item die kinematische Struktur der QM 
        (endlichdimensionaler symplektischer Raum),
  \item die kinematische Struktur aller freien Feldtheorien
        (unendlichdimensionale Einteilchenräume),
  \item sowie die formale Behandlung einfacher/interagierender QFT-Modelle
        (gleiche Algebra, veränderte Dynamik/Zustände)
\end{itemize}
einheitlich beschreibt.

Diese Ebene ist damit:
\begin{enumerate}
  \item \emph{allgemeiner} als die klassische Heisenberg-/Weyl-Algebra über \(\mathbb{R}^{2n}\),
  \item \emph{konkret genug}, um die vollständige Standard-QFT1/2-Abdeckung zu gewährleisten,
  \item \emph{kompatibel} mit dem noch allgemeineren C$^\ast$-/von-Neumann-Formalismus,
        in dem CCR/CAR-Algebren als strukturierte Spezialfälle erscheinen.
\end{enumerate}
\end{REM}
\end{unit}



\subsection{Endlich dimensionale QMs}

\begin{unit}
\emph{Endliche Quantenmechanik als Spezialfall der abstrakten CCR-Struktur}

In diesem Abschnitt zeigen wir explizit, wie die gewöhnliche nichtrelativistische
Quantenmechanik mit endlich vielen Freiheitsgraden in den abstrakten Rahmen
\[
(X,\sigma)\;\longmapsto\;\mathrm{CCR}(X,\sigma)
\]
eingebettet ist. Ziel ist es, die üblichen Objekte
\[
\mathcal{H}=L^2(\mathbb{R}^n),\quad
\hat{x}_j,\hat{p}_k,\quad
H=\tfrac{1}{2m}\hat{p}^2+V(\hat{x})
\]
als konkrete Realisierung der abstrakten Heisenberg- und CCR-Struktur zu identifizieren.

%--------------------------------------------------------------------
\emph{1. Kinematik: Der symplektische Einteilchenraum}

\emph{Abstrakt.}
Sei \(V\simeq\mathbb{R}^n\) der klassische Konfigurationsraum.
Der \emph{klassische Phasenraum} ist \(T^\ast V\simeq \mathbb{R}^{2n}\).
Wir setzen
\[
X := V_x \oplus V_p \;\cong\; \mathbb{R}^n_x\oplus\mathbb{R}^n_p\cong\mathbb{R}^{2n}
\]
mit der kanonischen symplektischen Form
\[
\sigma\bigl((x,p),(x',p')\bigr)
:= p\cdot x' - p'\cdot x.
\]
Damit ist \((X,\sigma)\) ein endlichdimensionaler symplektischer Raum.

\emph{Interpretation.}
Die Paare \((x,p)\in X\) repräsentieren klassische Zustände
\((\text{Ort},\text{Impuls})\), und \(\sigma\) ist die übliche symplektische Struktur
des Phasenraums.

%--------------------------------------------------------------------
\emph{2. Abstrakte Heisenberg-Algebra \texorpdfstring{$\mathfrak{h}(X)$}{h(X)}}

Nach der allgemeinen Definition (vgl.\ \textsc{MID-HEIS}) ist die
\emph{Heisenberg-Lie-Algebra}
\[
\mathfrak{h}(X) := X \oplus \mathbb{R}Z
\]
mit Lie-Klammer
\[
[(f,\alpha Z),(g,\beta Z)]
= (0,\ \sigma(f,g)\,Z).
\]
Wählt man eine Basis \(\{e_j\}_{j=1}^n\) von \(V_x\) und \(\{e^j\}_{j=1}^n\) von \(V_p\),
so kann man
\[
X_j := (e_j,0),\qquad
P_j := (0,e^j)
\]
als Erzeuger wählen. Dann lauten die Relationen
\[
[X_j,P_k] = \delta_{jk}\,Z,\qquad
[X_j,X_k]=[P_j,P_k]=[Z,\cdot]=0.
\]
Identifiziert man formal \(Z\sim i\hbar\), erhält man die bekannten
kanonischen Kommutationsrelationen:
\[
[X_j,P_k] = i\hbar\,\delta_{jk}.
\]

Die universelle Hüllalgebra \(U(\mathfrak{h}(X))\) enthält dann
Polynome in den \(X_j,P_k\), z.\,B. Hamiltonoperatoren der Form
\[
H = \frac{1}{2m}\sum_{j=1}^n P_j^2 + V(X_1,\dots,X_n).
\]

%--------------------------------------------------------------------
\emph{3. Die CCR-Weyl-C$^\ast$-Algebra \texorpdfstring{$\mathrm{CCR}(X,\sigma)$}{CCR(X,sigma)}}

\emph{Abstrakte Definition.}
Die CCR- (oder Weyl-) C$^\ast$-Algebra \(\mathrm{CCR}(X,\sigma)\)
ist die C$^\ast$-Vervollständigung der *-Algebra,
erzeugt von Symbolen \(W(f)\) (\(f\in X\)) mit den Weyl-Relationen
\[
W(f)W(g) = e^{-\frac{i}{2}\sigma(f,g)}W(f+g),\qquad
W(f)^\ast = W(-f).
\]

Diese Algebra ist ein abstrakter Vertreter der kanonischen
Kommutationsrelationen. Die \emph{infinitesimale} Heisenberg-Algebra
\(\mathfrak{h}(X)\) kann als Lie-Algebra der „logarithmischen Ableitungen“
der \(W(f)\) verstanden werden.

%--------------------------------------------------------------------
\emph{4. Schrödinger-Darstellung als konkrete Realisierung}

\emph{Hilbertraum.}
Wir wählen als Zustandsraum den Hilbertraum
\[
\mathcal{H} := L^2(\mathbb{R}^n,dx),
\]
mit kanonischem Skalarprodukt
\[
\langle \psi,\varphi\rangle = \int_{\mathbb{R}^n} \overline{\psi(x)}\varphi(x)\,dx.
\]

\emph{Ort- und Impulsoperatoren.}
Wir definieren (zunächst auf einem dichten Unterraum, etwa den Schwartzfunktionen)
\[
(\hat{X}_j\psi)(x) := x_j\psi(x),\qquad
(\hat{P}_j\psi)(x) := -\,i\hbar\,\frac{\partial}{\partial x_j}\psi(x).
\]
Dann gilt (formal)
\[
[\hat{X}_j,\hat{P}_k]\psi
= i\hbar\,\delta_{jk}\,\psi,\qquad
[\hat{X}_j,\hat{X}_k]=[\hat{P}_j,\hat{P}_k]=0,
\]
d.\,h. \(\hat{X}_j,\hat{P}_k\) realisieren die Relationen der Heisenberg-Algebra
\(\mathfrak{h}(X)\).

\emph{Weyl-Operatoren in der Schrödinger-Darstellung.}
Zu jedem \(f=(x,p)\in X\) definieren wir einen unitären Operator
\(\pi_{\mathrm{Sch}}(W(f))\) auf \(\mathcal{H}\) durch
\[
\bigl(\pi_{\mathrm{Sch}}(W(x,p))\psi\bigr)(y)
:= e^{\frac{i}{\hbar}\left(p\cdot y - \tfrac{1}{2}p\cdot x\right)}\,
   \psi(y-x).
\]

Man überprüft:
\begin{align*}
\pi_{\mathrm{Sch}}(W(x,p))\pi_{\mathrm{Sch}}(W(x',p'))
&= e^{-\frac{i}{2\hbar}\sigma((x,p),(x',p'))}\,
   \pi_{\mathrm{Sch}}(W(x+x',p+p')),\\
\pi_{\mathrm{Sch}}(W(x,p))^\ast
&= \pi_{\mathrm{Sch}}(W(-x,-p)),
\end{align*}
also erfüllt \(\pi_{\mathrm{Sch}}\) genau die Weyl-Relationen.
Damit ist
\[
\pi_{\mathrm{Sch}}:\mathrm{CCR}(X,\sigma)\longrightarrow \mathcal{B}(\mathcal{H})
\]
eine \emph{konkrete Darstellung} der abstrakten CCR-Algebra.

\emph{Infinitesimale Generatoren.}
Differenziert man \(t\mapsto \pi_{\mathrm{Sch}}(W(t e_j,0))\) und
\(t\mapsto \pi_{\mathrm{Sch}}(W(0,t e^j))\) bei \(t=0\), so erhält man
\[
\hat{X}_j = i\hbar\,\frac{d}{dt}\pi_{\mathrm{Sch}}(W(t e_j,0))\bigg|_{t=0},
\quad
\hat{P}_j = i\hbar\,\frac{d}{dt}\pi_{\mathrm{Sch}}(W(0,t e^j))\bigg|_{t=0}.
\]
Damit erscheinen die üblichen Ort- und Impulsoperatoren
als \emph{infinitesimale Generatoren} der Weyl-Darstellung
von \(\mathrm{CCR}(X,\sigma)\).

%--------------------------------------------------------------------
\emph{5. Stone--von Neumann und Eindeutigkeit (bis auf Äquivalenz)}

\begin{THEO}{SvN-1}{Stone--von Neumann (informell)}
Sei \(\pi:\mathrm{CCR}(X,\sigma)\to \mathcal{B}(\mathcal{H}_\pi)\)
eine irreduzible, \emph{regular} (d.\,h. stetige) Darstellung
der CCR-Algebra für endlichdimensionales \(X\).
Dann ist \(\pi\) zu der Schrödinger-Darstellung \(\pi_{\mathrm{Sch}}\)
unitär äquivalent.
\end{THEO}

\emph{Interpretation.}
Für eine gegebene endliche Anzahl von Freiheitsgraden \(n\) ist die
Struktur der Hilbertraumdarstellung der Heisenberg-/Weyl-Algebra
durch \(L^2(\mathbb{R}^n)\) (bis auf unitäre Äquivalenz) \emph{eindeutig}.
Damit ist der abstrakte Rahmen \((X,\sigma)\mapsto\mathrm{CCR}(X,\sigma)\)
für die endliche QM nicht nur realisierbar, sondern im Wesentlichen
\emph{vollständig}.

%--------------------------------------------------------------------
\emph{6. Hamiltonoperatoren und Dynamik in \texorpdfstring{$U(\mathfrak{h}(X))$}{U(h(X))}}

\emph{Hamiltonoperatoren als Elemente der Hüllalgebra.}
In der abstrakten Hüllalgebra \(U(\mathfrak{h}(X))\) können wir
Hamiltonoperatoren als Polynome in \(X_j,P_k\) definieren, z.\,B.
\[
H = \frac{1}{2m}\sum_{j=1}^n P_j^2 + V(X_1,\dots,X_n).
\]
In der Schrödinger-Darstellung wird daraus der übliche Operator
\[
\hat{H} = -\frac{\hbar^2}{2m}\Delta + V(x)
\]
auf \(L^2(\mathbb{R}^n)\).

\emph{Dynamik.}
Die Heisenberg-Zeitentwicklung eines Observablenoperators \(A\) wird
abstrakt durch die innere Derivation
\[
\dot{A}(t) = \frac{i}{\hbar}[H,A(t)]
\]
in \(U(\mathfrak{h}(X))\) beschrieben.
In der Schrödinger-Darstellung entspricht dies der üblichen Formel
\[
A_H(t) = e^{\frac{i}{\hbar}t\hat{H}} A_S e^{-\frac{i}{\hbar}t\hat{H}},
\]
während die Zustandsdynamik (Schrödinger-Bild) durch
\(i\hbar\partial_t\psi = \hat{H}\psi\) gegeben ist.

Damit ist die volle Dynamik der nichtrelativistischen QM
ein Spezialfall der abstrakten Dynamik in \(U(\mathfrak{h}(X))\)
und ihrer Darstellung auf \(L^2(\mathbb{R}^n)\).

%--------------------------------------------------------------------
\emph{7. Fazit: Endliche QM als Spezialfall der Zwischenebene}

\begin{REM}{QM-SPECIAL}{Zusammenfassung}
Die gewöhnliche Quantenmechanik mit endlich vielen Freiheitsgraden
entsteht aus der Zwischenebene
\[
(X,\sigma)\longmapsto \mathrm{CCR}(X,\sigma)
\]
spezifisch durch die Wahl:
\begin{itemize}
  \item \(X = \mathbb{R}^{2n}\) mit kanonischer symplektischer Form \(\sigma\).
  \item CCR-Weyl-C$^\ast$-Algebra \(\mathcal{A} = \mathrm{CCR}(X,\sigma)\) als Observablenalgebra.
  \item Schrödinger-Darstellung \(\pi_{\mathrm{Sch}}:\mathcal{A}\to\mathcal{B}(L^2(\mathbb{R}^n))\),
        die nach Stone--von Neumann (unter Regularität) bis auf unitäre Äquivalenz eindeutig ist.
  \item Heisenberg-Algebra \(\mathfrak{h}(X)\) und Hüllalgebra \(U(\mathfrak{h}(X))\),
        deren Generatoren sich als die üblichen Ort- und Impulsoperatoren
        realisieren.
  \item Hamiltonoperatoren \(H \in U(\mathfrak{h}(X))\) und die daraus
        abgeleitete Dynamik \(\dot{A}=\frac{i}{\hbar}[H,A]\) im Heisenberg-Bild
        bzw.\ \(i\hbar\partial_t\psi=\hat{H}\psi\) im Schrödinger-Bild.
\end{itemize}

In diesem präzisen Sinn ist die endliche nichtrelativistische QM
ein \emph{Spezialfall} der abstrakten Heisenberg-/CCR-Struktur:
die allgemeineren Konstruktionen (symplektischer Einteilchenraum,
Heisenberg-Algebra, CCR-Algebra, Automorphismengruppe)
reduzieren sich für \(X=\mathbb{R}^{2n}\) genau auf das übliche
Quantensystem \((L^2(\mathbb{R}^n),\hat{x},\hat{p},\hat{H})\).
\end{REM}




\emph{Ergänzung 1: Abstrakte vs.\ konkrete Operatoren, Heisenberg- und Schrödinger-Bild}

\emph{(a) Abstrakte vs.\ konkrete Ebene.}

Im abstrakten CCR-Rahmen unterscheidet man streng zwischen

\begin{itemize}
  \item der \emph{abstrakten Observablenalgebra}
  \[
    \mathcal{A} := \mathrm{CCR}(X,\sigma),
  \]
  einer C$^\ast$-Algebra, deren Elemente wir formal als
  \emph{Observablen} (bzw.\ deren „exponentialisierte“ Versionen) interpretieren,
  und

  \item einer \emph{konkreten Darstellung}
  \[
    \pi : \mathcal{A} \longrightarrow \mathcal{B}(\mathcal{H}),
  \]
  die jedem abstrakten Observablen-Element \(A\in \mathcal{A}\) 
  einen \emph{beschränkten Operator} \(\pi(A)\) auf einem Hilbertraum \(\mathcal{H}\) zuordnet.
\end{itemize}

Die abstrakte Ebene \(\mathcal{A}\) ist „repräsentationsfrei“;
erst durch \(\pi\) werden die Algebraelemente als konkrete Operatoren auf Zuständen 
(konkret: Vektoren \(\psi\in\mathcal{H}\)) realisiert.

\emph{Wichtig:}  
\(\mathcal{A}\) ist \emph{kein} Unterraum von \(\mathcal{B}(\mathcal{H})\) per se, 
sondern ein separates Objekt; erst die Darstellung \(\pi\) stellt eine Brücke her:
\[
A\in\mathcal{A}\quad\mapsto\quad \pi(A)\in\mathcal{B}(\mathcal{H}).
\]

\medskip

\emph{(b) Abstrakte Dynamik vs.\ konkrete Dynamik.}

Auf der abstrakten Ebene  
ist eine \emph{Dynamik} gegeben durch eine einparametrige Gruppe von
$\ast$-Automorphismen
\[
  \alpha_t:\mathcal{A}\to\mathcal{A},\qquad
  \alpha_{t+s} = \alpha_t\circ\alpha_s,\quad\alpha_0=\mathrm{id}.
\]
Man schreibt (formal)
\[
  \frac{d}{dt}\,\alpha_t(A) = \delta(A_t),\qquad \delta:\mathcal{A}\to\mathcal{A}
\]
für eine (oft innere) Derivation \(\delta\).

In einer Darstellung \(\pi\) auf \(\mathcal{H}\) ist diese abstrakte Dynamik
durch eine strongly continuous unitäre Gruppe
\[
  U(t):\mathcal{H}\to\mathcal{H},\qquad U(t)U(s)=U(t+s),
\]
realisiert, so dass
\[
  \pi(\alpha_t(A)) = U(t)\,\pi(A)\,U(t)^{-1}.
\]

\medskip

\emph{(c) Heisenberg- und Schrödinger-Bild präzise getrennt.}

\textbf{Heisenberg-Bild (algebraischer Standpunkt).}
\begin{itemize}
  \item Auf der abstrakten Ebene: wir betrachten \(\alpha_t(A)\in\mathcal{A}\), d.\,h.\ 
        die Observablen selbst sind zeitabhängig.
  \item In einer Darstellung \(\pi\): 
        \[
          A_H(t) := \pi(\alpha_t(A)) = U(t)\pi(A)U(t)^{-1}.
        \]
        Zustände bleiben (formal) zeitlich konstant.
\end{itemize}

\textbf{Schrödinger-Bild (zustandsorientierter Standpunkt).}
\begin{itemize}
  \item Man fixiert die Observablen \(\pi(A)\) als zeitunabhängig.
  \item Die Zeitentwicklung wird stattdessen auf Zustände verschoben:
        \[
          \psi_S(t) := U(t)^{-1}\psi(0),
        \]
        und Erwartungswerte werden durch
        \[
          \langle \psi_S(t), \pi(A)\psi_S(t)\rangle
        \]
        gegeben.
\end{itemize}

Formal sind beide Bilder äquivalent, da
\[
  \langle \psi, A_H(t)\psi\rangle
  \;=\;
  \langle \psi_S(t), \pi(A)\psi_S(t)\rangle
\]
für alle \(A\in\mathcal{A}\) und initialen Zustände \(\psi\) gilt.

\medskip

\emph{(d) Zustände vs.\ Operatoren.}

\begin{itemize}
  \item Ein \emph{(algebraischer) Zustand} auf \(\mathcal{A}\)
        ist eine positive normierte Linearform
        \(\omega:\mathcal{A}\to\mathbb{C}\) mit \(\omega(\mathbf{1})=1\).
  \item In einer Darstellung \(\pi\) mit Hilbertraum \(\mathcal{H}\)
        entspricht ein Vektorzustand \(\omega_\psi(A):=\langle \psi,\pi(A)\psi\rangle\)
        einem \emph{Zustandsvektor} \(\psi\in\mathcal{H}\).
\end{itemize}

Damit ist klar:
\[
\boxed{
\text{Zustände leben (repräsentiert) als Vektoren } \psi\in\mathcal{H},\;
\text{Observablen als Operatoren } \pi(A)\in\mathcal{B}(\mathcal{H}).
}
\]

Die abstrakte CCR-/Weyl-Algebra \(\mathcal{A}\) kodiert die Struktur der Observablen,
die Darstellung \(\pi\) und ein Zustand \(\psi\) kodieren die konkrete Physik
(Erwartungswerte, Spektren, Dynamik).


%====================================================================
\emph{Ergänzung 2: CCR als Ursprung der Observablen und ihre Generatorenrolle}

\emph{(a) CCR als abstrakte Quelle der üblichen Operatoren.}

In vielen QM-Vorlesungen erscheinen die Orts- und Impulsoperatoren
\(\hat{x}_j,\hat{p}_k\) zunächst „ad hoc“ als
\[
(\hat{x}_j\psi)(x) = x_j\psi(x),\qquad
(\hat{p}_j\psi)(x) = -i\hbar\partial_{x_j}\psi(x).
\]

Im CCR-Rahmen sieht man diese Operatoren als \emph{konkrete Realisierung}
von abstrakten Strukturelementen:
\begin{itemize}
  \item Man startet mit dem symplektischen Raum \((X,\sigma)\).
  \item Bildet daraus die abstrakte CCR-C$^\ast$-Algebra \(\mathcal{A}=\mathrm{CCR}(X,\sigma)\)
        mit Weyl-Generatoren \(W(f)\), \(f\in X\).
  \item Wählt dann eine Darstellung \(\pi:\mathcal{A}\to\mathcal{B}(\mathcal{H})\).
\end{itemize}

Erst in dieser Darstellung erscheinen die „bekannten“ Operatoren
\(\hat{x}_j,\hat{p}_k\) als abgeleitete Objekte,
nämlich als \emph{infinitesimale Generatoren} bestimmter unitärer Einparametergruppen
\(t\mapsto \pi(W(tf))\).

Dies macht klar:
\[
\boxed{
\text{„Ort“ und „Impuls“ sind nicht primäre Operatoren, 
sondern Generatoren der abstrakten CCR-Struktur.}
}
\]

\medskip

\emph{(b) Weyl-Relationen und Generatoren (Stone-Theorem).}

Für \(f\in X\) liefert \(t\mapsto W(tf)\) in \(\mathcal{A}\)
eine abstrakte Einparametergruppe.
In einer (regularen) Darstellung \(\pi\) wird daraus
eine strongly continuous unitäre Gruppe
\[
U_f(t) := \pi(W(tf)) \in \mathcal{B}(\mathcal{H}).
\]

Nach dem Satz von Stone existiert ein (wesentlich) selbstadjungierter Operator
\(A_f\) auf \(\mathcal{H}\) mit
\[
U_f(t) = e^{itA_f}.
\]
Man nennt \(A_f\) den \emph{(infinitesimalen) Generator} der Gruppe \(U_f(t)\).

\emph{Beispiele.}
\begin{itemize}
  \item Wählt man \(f\) so, dass \(W(tf)\) räumliche Translationen generiert,
        dann ist der Generator der Impulsoperator \(\hat{p}\).
  \item Wählt man \(f\) so, dass \(W(tf)\) Phasenmodulationen generiert,
        dann ist der Generator der Ortsoperator \(\hat{x}\).
\end{itemize}

Damit werden Observablen als „Differentialquotienten“ von physikalisch
bedeutenden Symmetrien (Translationsgruppen, Phasenrotationen, etc.)
sichtbar:
\[
\hat{A}_f = \left.\frac{1}{i}\frac{d}{dt}U_f(t)\right|_{t=0}.
\]

\medskip

\emph{(c) Konsequenzen: Symmetrien, Spektraltheorie, Messinterpretation.}

\begin{itemize}
  \item \textbf{Symmetrien als unitäre Gruppen:}
        Jede physikalisch relevante kontinuierliche Symmetrie
        wird durch eine unitäre Gruppe \(U(t)\) auf \(\mathcal{H}\) beschrieben.
        Der zugehörige selbstadjungierte Generator \(A\) ist die „zugehörige Observable“
        (z.\,B.\ Energie, Impuls, Drehimpuls).

  \item \textbf{Spectrum = mögliche Messergebnisse:}
        Die Spektralzerlegung eines selbstadjungierten Operators \(A\)
        gibt die möglichen Messergebnisse dieser Observable wieder.
        Die CCR-/Weyl-Struktur legt fest, wie diese Spektren
        zueinander in Relation stehen (z.\,B.\ Unschärferelationen).

  \item \textbf{Dynamik als Generator:}
        Der Hamiltonoperator \(\hat{H}\) ist der Generator der Zeittranslationsgruppe
        \(U(t)=e^{-i t\hat{H}/\hbar}\).
        Damit ist die Dynamik konsequent als weitere Einparametergruppe
        im selben abstrakten CCR-/C$^\ast$-Rahmen beschrieben.
\end{itemize}

\medskip

\emph{(d) Zusammenfassung der „Generator-Perspektive“.}

\[
\boxed{
\begin{aligned}
&\text{Abstrakt: } (X,\sigma)\ \leadsto\ \mathrm{CCR}(X,\sigma)
\text{ mit Weyl-Generatoren }W(f).\\[4pt]
&\text{Darstellung: } \pi(W(f)) \text{ ergibt unitäre Gruppen }U_f(t),\\
&\text{deren Generatoren }A_f \text{ als Observablen interpretiert werden.}\\[4pt]
&\Rightarrow\ \text{„Ort“, „Impuls“, „Energie“ usw.\ sind keine ad hoc-Operatoren,}\\
&\quad\text{sondern infinitesimale Generatoren von Symmetrien, die bereits in den CCR stecken.}
\end{aligned}
}
\]

In diesem Sinn liefert der abstrakte CCR-Rahmen eine \emph{konzeptionelle Rechtfertigung}
der im Vorlesungsalltag oft „einfach eingeführten“ Operatoren:
sie sind nicht zufällig gewählt, sondern entstehen systematisch als Generatoren
von unitären Darstellungen der zugrundeliegenden symplektischen Struktur \((X,\sigma)\)
und ihrer Symmetrien.

\end{unit}



\begin{unit}
\emph{Einbettung des Lie--Gruppen- und Flussbildes in den CCR-Rahmen}

In diesem Abschnitt wird präzisiert, wie die zuvor entwickelte
\emph{Heisenberg/CCR-/Weyl-Struktur} mit dem \emph{Lie-Gruppen- und Flussbild}
(Zeitentwicklung, Symmetrien, Orbits) zusammenhängt.
Grobes Bild:

\[
\boxed{
\text{CCR-/Heisenberg-Algebra} \;\;+\;\; \text{unitäre Darstellung von $G$}
\;\Longrightarrow\; \text{Flüsse von Zuständen \& Observablen}
}
\]

Wir ordnen dazu die Bausteine deines Lie-Gruppen-Textes in den CCR-Rahmen ein
und machen die Verknüpfungen explizit.

%---------------------------------------------------------------
\emph{Horizontale Achse: Kinematik über CCR und Heisenberg-Algebra}

\emph{Abstrakte kinematische Seite.}

\begin{itemize}
  \item Ausgangspunkt ist ein symplektischer Raum \((X,\sigma)\)
        (klassischer Phasenraum / Einteilchenraum).
  \item Daraus konstruieren wir
        \begin{itemize}
          \item die \emph{Heisenberg-Lie-Algebra} \(\mathfrak{h}(X) = X\oplus\R Z\),
          \item ihre Hüllalgebra \(U(\mathfrak{h}(X))\),
          \item und die CCR-/Weyl-C$^\ast$-Algebra
                \(\mathcal{A} := \mathrm{CCR}(X,\sigma)\).
        \end{itemize}
  \item \(\mathcal{A}\) ist die \emph{abstrakte Observablenalgebra}:
        ihre (selbstadjungierten) Elemente modellieren mögliche Observablen.
\end{itemize}

\emph{Konkrete kinematische Realisierung.}

\begin{itemize}
  \item Eine Schrödinger-Darstellung
        \(\pi:\mathcal{A}\to\mathcal{B}(\mathcal{H})\) auf einem Hilbertraum
        \(\mathcal{H}=L^2(\R^n)\) realisiert die abstrakten \(W(f)\) als konkrete
        Operatoren \(\pi(W(f))\) (Translation, Phasenrotation usw.).
  \item Die infinitesimalen Generatoren der Weyl-Operatoren sind die üblichen
        \(\hat x_j,\hat p_k\), die wir aus der Vorlesung kennen; sie leben als
        unbeschränkte, selbstadjungierte Operatoren in \(\mathcal{B}(\mathcal{H})\).
\end{itemize}

Bis hierher wird \emph{nur} festgelegt, welche Observablen existieren und wie
sie kommutieren: reine Kinematik.

%---------------------------------------------------------------
\emph{Vertikale Achse: Lie-Gruppen, Flüsse und Zustandsentwicklung}

Jetzt kommt der Teil ins Spiel, den du in deinem Abschnitt
\emph{„Lie-Gruppen, Flüsse und die Entwicklung von Zustandsbeschreibungen“}
formuliert hast. Dort sind die Grundobjekte:

\begin{itemize}
  \item Eine Lie-Gruppe \(G\) mit Lie-Algebra \(\mathfrak g\).
  \item Eine unitäre Darstellung
        \[
          U:G\to\mathcal{U}(\mathcal{H}),
        \]
        die jeder Symmetrie (oder Zeitentwicklung) ein unitäres \(U(g)\) zuordnet.
  \item Flüsse \(t\mapsto \exp(tX)\) in \(G\) (Einparameteruntergruppen)
        und induzierte Flüsse
        \(\Phi_X(t,p)=\exp(tX)\cdot p\) auf einem Zustandsraum \(M\).
\end{itemize}

Im Quantenkontext setzen wir dabei:
\[
M := \mathbb{P}(\mathcal{H})
\quad\text{(projektiver Hilbertraum, Raum der Strahlen)}.
\]
Ein Fluss in \(G\) erzeugt via \(U\) einen Fluss von Zustandsbeschreibungen
\(\psi(t)=U(\exp tX)\psi_0\) und damit einen Fluss von physikalischen Zuständen
\([\,\psi(t)\,]\in\mathbb{P}(\mathcal{H})\).

\emph{Dein Text zu Lie-Gruppenwirkungen, Flüssen und Orbits
beschreibt genau diese \emph{vertikale} Ebene:  
wie Symmetriegruppen und Zeitentwicklung auf den Zustand(sbeschreibungen)
wirken.}

%---------------------------------------------------------------
\emph{Brücke 1: Von CCR-Automorphismen zu $G$-Flüssen auf Zuständen}

\emph{Abstrakte Symmetrie auf der Observablen-Ebene.}

Sei \(G\) nun eine Symmetriegruppe (z.\,B.\ Translationen, Rotationen,
Zeittranslationen), die \emph{abstrakt} auf der CCR-Algebra \(\mathcal{A}\)
agiert:
\[
\alpha: G\to\mathrm{Aut}(\mathcal{A}),
\qquad
g\mapsto\alpha_g,
\]
wobei \(\alpha_g\) ein $\ast$-Automorphismus ist (Observablen werden zu
Observablen, Spektrum bleibt erhalten etc.).

Beispiel:
\[
\alpha_g(W(f)) = W(S_g f),\qquad S_g\in\mathrm{Sp}(X,\sigma)
\]
für eine symplektische Transformation \(S_g\) des Einteilchenraums.

\emph{Implementierung durch unitäre Darstellung.}

In einer Darstellung \(\pi:\mathcal{A}\to\mathcal{B}(\mathcal{H})\)
sollen die abstrakten Symmetrien \(\alpha_g\) durch unitäre Operatoren
\(U(g)\) implementiert werden:
\[
\pi(\alpha_g(A)) = U(g)\,\pi(A)\,U(g)^{-1},
\qquad A\in\mathcal{A}.
\]
Dies ist genau die Kovarianzbedingung, die du in deinem Lie-Gruppen-Teil
als
\[
dU(\mathrm{Ad}_g X)=\mathrm{Ad}_{U(g)}(dU(X))
\]
auf Generatorniveau formuliert hast.

\emph{Zustandsflüsse.}

Auf Zuständen (Vektoren) definiert \(U(g)\) dann den Fluss
\[
\psi\mapsto U(g)\psi,\qquad
[\psi]\mapsto [U(g)\psi],
\]
so dass Erwartungswerte kovariant sind:
\[
\langle \psi,\pi(A)\psi\rangle
=
\langle U(g)\psi,\pi(\alpha_g(A))U(g)\psi\rangle.
\]

Damit ist der Abschnitt über
\emph{„Lie-Gruppenwirkungen, Flüsse und die Entwicklung von Zustandsbeschreibungen“}
genau die \emph{konkrete realisierte} Form der abstrakten Symmetrie-Automorphismen
\(\alpha_g\) der CCR-Algebra.

%---------------------------------------------------------------
\emph{Brücke 2: Generatoren als Observablen \& CCR als Herkunft}

Deine Lie-Gruppen-Abschnitte arbeiten intensiv mit Generatoren:

\begin{itemize}
  \item Für eine Einparameteruntergruppe
        \(t\mapsto g(t)=\exp(tX)\in G\) gilt
        \[
          U(g(t)) = e^{t\,dU(X)} = e^{-\frac{i}{\hbar}t\hat{X}}.
        \]
  \item Die \(\hat{X}=i\hbar dU(X)\) sind \emph{selbstadjungierte Observablen}
        (Energie, Impuls, Drehimpuls, Ladung, \dots).
  \item Dein „Flows on $M$“–Teil beschreibt:
        \[
          \psi(t)=U(g(t))\psi_0,\qquad
          [\psi(t)]\in\mathbb{P}(\mathcal{H})
        \]
        als Entwicklung von Zustandsbeschreibungen.
\end{itemize}

Im CCR-Rahmen sehen wir jetzt:

\begin{itemize}
  \item Die abstrakte Heisenberg-/CCR-Algebra bestimmt, welche
        Kommutatoren \([\hat X,\hat P]\), \([\hat J_i,\hat J_j]\), \dots
        zulässig sind (CCR-Relationen, Lie-Algebra-Relationen).
  \item Die Lie-Gruppen-Darstellung \(U:G\to\mathcal{U}(\mathcal{H})\)
        ist nichts anderes als die exponentielle Integration der Generatoren,
        die \emph{Elemente der CCR-/Heisenberg-/Lie-Algebra} sind.
  \item Deine Flussgleichungen
        \[
          \dot g(t)=(L_{g(t)})_*X,\quad
          \dot\psi(t)=-\frac{i}{\hbar}\hat{H}\psi(t)
        \]
        sind genau die global integrierte Form dessen, was wir auf
        \emph{algebraischer Seite} als
        \(\dot{A}(t)=\frac{i}{\hbar}[H,A(t)]\)
        und \(\dot{\omega}_t(\cdot)=\omega_t(i[H,\cdot])\) schreiben.
\end{itemize}

Kurz:
\[
\boxed{
\text{Generatoren } \hat X
\text{ sind Observablen in der CCR-Algebra,}
\quad
\text{Lie-Flüsse sind ihre exponentiellen Einparametergruppen.}
}
\]

%---------------------------------------------------------------
\emph{Brücke 3: Deine „Flüsse in der Lie-Gruppe“ als universeller Oberbau}

Dein Abschnitt

\begin{quote}
\emph{„Flüsse in der Lie-Gruppe und induzierte Flüsse“} \\
\emph{„Zeitunabhängige vs. zeitabhängige Generatoren“} \\
\emph{„Invarianz, Orbits und Erhaltung“}
\end{quote}

passt in den CCR-Rahmen wie folgt:

\emph{(i) Zeitunabhängige und zeitabhängige Generatoren.}

\begin{itemize}
  \item Zeitunabhängiger Generator \(X\in\mathfrak g\):  
        \(g(t)=\exp(tX)\), \(U(t)=U(\exp tX)\), Heisenberg-Bild:
        \[
        \psi(t)=U(t)\psi_0
        \quad\text{oder}\quad
        A_H(t)=U(t)^{-1} A_S U(t).
        \]
        Das ist genau das, was im CCR-Bild durch einen festen Hamiltonoperator \(H\in U(\mathfrak{h}(X))\)
        beschrieben wird.
  \item Zeitabhängiger Generator \(X(t)\):  
        \(g(t)=\mathcal{P}\exp\!\int_0^t X(s)\,ds\).
        Im CCR-/QM-Setting ist dies ein zeitabhängiger Hamiltonoperator
        \(\hat H(t)\) und eine zeitgeordnete Entwicklung
        \[
          U(t)=\mathcal{P}\exp\!\Big(-\frac{i}{\hbar}\int_0^t\hat H(s)\,ds\Big).
        \]
        Hier sind deine \(\mathcal{P}\)-Exponentialabbildungen direkt die
        in der Physik übliche zeitgeordnete Exponentialentwicklung (Dyson-Reihen).
\end{itemize}

\emph{(ii) Invarianz, Orbits, Erhaltung.}

\begin{itemize}
  \item Invariante Funktion \(f\) auf \(M\):  
        \(f(g\cdot p)=f(p)\).  
        Im QM/CCR-Setting: Erwartungswertfunktion
        \(\psi\mapsto\langle\psi,\hat A\psi\rangle\) ist invariant unter \(U(G)\)
        genau dann, wenn
        \[
          U(g)\hat A U(g)^{-1}=\hat A,\quad \forall g\in G,
        \]
        bzw.\ infinitesimal \([\hat A, dU(\mathfrak g)]=0\).
  \item Orbits \(\mathcal{O}_\psi = \{U(g)\psi\mid g\in G\}\):  
        Alle Beschreibungen desselben physikalischen Zustands (bis zu Symmetrien).
        Das passt zur Interpretation der CCR-Algebra:
        dieselbe physikalische Observable \(\hat A\) kann in verschiedenen
        Darstellungen (z.\,B.\ verschiedenen Beobachterrahmen) unterschiedlich aussehen,
        bleibt aber durch die Adjoint-Aktion \(U(g)\hat A U(g)^{-1}\) verbunden.
  \item Casimirs: In deinem Text als invariantes \(f\) oder zentraler Operator;  
        im CCR-/Lie-Algebra-Rahmen sind Casimir-Operatoren zentrale Elemente,
        deren Eigenwerte \emph{Quantenzahlen} (Spin, Masse, Ladung) klassifizieren.
\end{itemize}

%---------------------------------------------------------------
\emph{Brücke 4: Integration in die Skript-Struktur}

Didaktisch kannst du die Integration etwa so anordnen:

\begin{enumerate}
  \item \textbf{Kinematische Ebene (Heisenberg/CCR):}
        \begin{itemize}
          \item Definition von \((X,\sigma)\), \(\mathfrak{h}(X)\), \(U(\mathfrak{h}(X))\),
                \(\mathrm{CCR}(X,\sigma)\).
          \item Schrödinger-Darstellung als Spezialfall.
        \end{itemize}
  \item \textbf{Dynamik und Hamiltonoperatoren:}
        \begin{itemize}
          \item Hamiltonoperator \(H\in U(\mathfrak{h}(X))\),
                Heisenberg-Gleichung \( \dot A = \frac{i}{\hbar}[H,A] \),
                Schrödinger-Gleichung \( i\hbar\dot\psi=\hat H\psi \).
        \end{itemize}
  \item \textbf{Symmetrien als Lie-Gruppenaktionen:}
        \begin{itemize}
          \item Hier fügst du \emph{deinen gesamten Lie-Gruppen-Text} ein:
            \begin{itemize}
              \item Grundobjekte \(G,\mathfrak g,\mathcal{H}\).
              \item Flüsse in \(G\) und induzierte Flüsse auf \(\mathbb{P}(\mathcal{H})\).
              \item Zeitunabhängige vs.\ zeitabhängige Generatoren (\(\mathcal{P}\exp\)).
              \item Invarianz, Orbits, Noether-Idee, Casimirs.
            \end{itemize}
          \item Parallel deutest du an, wie \(G\) als Symmetriegruppe
                von \(\mathcal{A}=\mathrm{CCR}(X,\sigma)\) verweist:
                \(\alpha_g\in\mathrm{Aut}(\mathcal{A})\), implementiert durch \(U(g)\).
        \end{itemize}
  \item \textbf{Zusammenführung:}
        \begin{itemize}
          \item Diagramm:
            \[
            \begin{tikzcd}[column sep=huge,row sep=large]
            (X,\sigma) \arrow{r}{\text{CCR}} \arrow{d}{S_g}
              & \mathcal{A}=\mathrm{CCR}(X,\sigma)
                \arrow{d}{\alpha_g} \arrow{r}{\pi}
              & \pi(\mathcal{A})\subset\mathcal{B}(\mathcal{H}) \arrow{d}{\mathrm{Ad}_{U(g)}}\\
            (X,\sigma) \arrow{r}{\text{CCR}}
              & \mathcal{A} \arrow{r}{\pi}
              & \pi(\mathcal{A})
            \end{tikzcd}
            \]
          \item Lie-Flüsse in \(G\) correspondieren zu unitären Flüssen in \(\mathcal{U}(\mathcal{H})\),
                die wiederum Flüsse von Zustandsbeschreibungen und Observablen induzieren.
        \end{itemize}
\end{enumerate}

%---------------------------------------------------------------
\emph{Bemerkung zur Rollenverteilung}

\begin{REM}{Lie-CCR-Integration}{Rollenverteilung}
Die beiden Blöcke \emph{CCR/Heisenberg} und \emph{Lie-Gruppen/Flüsse} 
spielen unterschiedliche, aber komplementäre Rollen:

\begin{itemize}
  \item Die CCR-/Heisenberg-Struktur beschreibt die \emph{interne Nichtkommutativität}
        der Observablen (Ort, Impuls, Spin \dots), also die kinematische Quantenstruktur.
  \item Die Lie-Gruppen-/Fluss-Struktur beschreibt die \emph{äußeren Transformationen}
        (Symmetrien, Zeitentwicklung) als Flüsse in \(G\) und \(\mathcal{U}(\mathcal{H})\),
        und wie sie Zustände und Observablen bewegen.
  \item Die Verbindung besteht darin, dass die Generatoren der Lie-Gruppenflüsse
        selbst Elemente der Heisenberg-/CCR-/Lie-Algebren sind:  
        sie sind \emph{Observablen}, deren Flüsse wir als Prozesse (Dynamik,
        Symmetrieoperationen) interpretieren.
\end{itemize}

Integriert man diese Sichtweisen, erhält man ein sehr einheitliches Bild:
\[
\boxed{
\begin{aligned}
&\text{CCR/Heisenberg beschreiben, \emph{was} beobachtbar ist (Struktur der Operatoren).}\\
&\text{Lie-Gruppen und Flüsse beschreiben, \emph{wie} sich diese Beobachtungen transformieren.}
\end{aligned}
}
\]
\end{REM}

\end{unit}



\begin{unit}
\emph{Top--Down-Architektur: Von Physik zu CCR, Lie-Gruppen und Flüssen}

Ziel dieses Abschnitts ist eine \emph{konzeptionelle} Herleitung der verwendeten
Objekte und ihrer Funktionen, mit explizitem Blick auf den diagrammatischen Zusammenhang.

Ganz grob:

\[
\boxed{
\text{Physikalische Wünsche (Zustände, Observablen, Symmetrien, Dynamik)}
\;\leadsto\;
\text{Mathematische Objekte (CCR, Hilberträume, Lie-Gruppen, Flüsse)}
}
\]

Wir gehen in Schichten vor.

%--------------------------------------------------------------------
\emph{Schicht 0: Physikalische Anforderungen}

\emph{Ausgangspunkt:} Welche Struktur soll eine Theorie \emph{mindestens} liefern?

\begin{itemize}
  \item[(a)] \textbf{Zustände:} 
  Es soll eine Menge \(\mathcal{S}\) von Zuständen geben (z.\,B. klassische
  Phasenraumspunkte, quantenmechanische Zustände), die das System beschreiben.

  \item[(b)] \textbf{Observablen:}
  Es soll eine Menge \(\mathcal{O}\) messbarer Größen geben (Ort, Impuls, Energie, \dots),
  die reelle Messwerte liefern und auf Zuständen ausgewertet werden können.
  Erwartungswerte sollen sinnvoll definiert sein.

  \item[(c)] \textbf{Symmetrien:}
  Es gibt Transformationen \(g\) (Beobachterwechsel, Rotationen, Translationen, Eichungen),
  die Zustände und Observablen verändern, aber physikalische Aussagen invariant lassen.

  \item[(d)] \textbf{Dynamik (Zeitentwicklung):}
  Es gibt eine ausgezeichnete „Zeit“-Transformation (oder Gruppe \(\mathbb{R}\)),
  die beschreibt, wie sich Zustände und Observablen zeitlich entwickeln.
\end{itemize}

Wir wollen Objekte konstruieren, die \emph{genau diese Funktionen} mathematisch realisieren.

%--------------------------------------------------------------------
\emph{Schicht 1: Klassische Kinematik -- symplektischer Phasenraum}

\emph{Funktion dieser Schicht:} 
Kodierung von \emph{Zuständen und Observablen} in der klassischen Mechanik,
ohne Quantisierung und ohne explizite Symmetrien.

\begin{itemize}
  \item Der \textbf{Zustandsraum} ist ein glatter Phasenraum
        \[
          M \cong T^\ast V \cong \R^{2n},
        \]
        versehen mit einer symplektischen Form \(\omega\).

  \item \textbf{Observablen} sind glatte Funktionen
        \[
          f:M\to\R.
        \]
        Sie bilden eine Poisson-Algebra
        \((C^\infty(M),\{\cdot,\cdot\})\), mit
        \[
          \{f,g\}=\omega(X_f,X_g),
        \]
        und \(X_f\) dem hamiltonschen Vektorfeld.

  \item \textbf{Kinematik}:
        Die Struktur \((M,\omega)\) kodiert, wie Observablen \emph{zueinander stehen};
        Dynamik entsteht erst mit der Wahl einer speziellen Funktion \(H\) (Hamiltonfunktion).
\end{itemize}

In der endlichdimensionalen QM-Vorlesung wählen wir typischerweise
\[
X:=V_x\oplus V_p \cong \R^{2n}
\]
mit kanonischer symplektischer Form \(\sigma\).
Damit spielen \((X,\sigma)\) die Rolle der \emph{linearen Version} des klassischen Phasenraums.

%--------------------------------------------------------------------
\emph{Schicht 2: Quantenkinematik -- Heisenberg- und CCR-Algebra}

\emph{Funktion dieser Schicht:}
Kodierung der \emph{quantenmechanischen Kinematik}:  
nichtkommutative Observablen, Unschärferelationen, Operatorstruktur.

\emph{(1) Heisenberg-Lie-Algebra \(\mathfrak{h}(X)\).}

Aus \((X,\sigma)\) wird zunächst eine Lie-Algebra konstruiert:

\[
\mathfrak{h}(X) := X \oplus \R Z,
\qquad
[(f,\alpha Z),(g,\beta Z)] = (0,\sigma(f,g)Z).
\]

In Koordinaten:
\[
[X_j,P_k] = i\hbar\,\delta_{jk},\qquad [X_j,X_k]=[P_j,P_k]=[Z,\cdot]=0.
\]

\begin{itemize}
  \item \(\mathfrak{h}(X)\) ist die \emph{infinitesimale} Kodierung der 
        kanonischen Kommutationsrelationen: sie ersetzt die klassische Poisson-Klammer
        durch einen Kommutator.
  \item Sie ist rein \emph{kinematisch}: sie sagt nur, wie „Ort“ und „Impuls“ 
        zueinander stehen, nicht, wie sie sich in der Zeit entwickeln.
\end{itemize}

\emph{(2) Hüllalgebra \(U(\mathfrak{h}(X))\).}

Die universelle Hüllalgebra
\[
U(\mathfrak{h}(X))
= T(\mathfrak{h}(X)) / \langle xy-yx-[x,y]\rangle
\]
ist die kleinste assoziative Algebra, die \(\mathfrak{h}(X)\) und ihre Klammer enthält.

\begin{itemize}
  \item Sie enthält alle \emph{Operatorpolynome in \(X_j,P_k\)}:
        \[
        H = \frac{1}{2m} P^2 + V(X)
        \]
        lebt \emph{abstrakt} in \(U(\mathfrak{h}(X))\).
  \item Diese Ebene beschreibt die \emph{Dynamik der Observablen} als innere Derivationen:
        \(\dot A = \frac{i}{\hbar}[H,A]\).
\end{itemize}

\emph{(3) CCR-/Weyl-C$^\ast$-Algebra \(\mathrm{CCR}(X,\sigma)\).}

Um eine C$^\ast$-Struktur und beschränkte Operatoren zu bekommen, exponentiieren wir
die Heisenberg-Algebra:

\[
\mathrm{CCR}(X,\sigma)
:= C^\ast\bigl(\,W(f)\mid f\in X\ \text{mit Weyl-Relationen}\,\bigr),
\]
\[
W(f)W(g)=e^{-\frac{i}{2}\sigma(f,g)}W(f+g),\quad W(f)^\ast=W(-f).
\]

\begin{itemize}
  \item \(\mathrm{CCR}(X,\sigma)\) ist die \emph{abstrakte Observablenalgebra}:
    selbstadjungierte Elemente \(A=A^\ast\) sind abstrakte Observablen.
  \item Sie verallgemeinert die Heisenberg-Algebra auf beschränkte, 
        exponentialisierte Operatoren.
\end{itemize}

\medskip

\noindent
\textbf{Diagramm (kinematisch):}

\[
\begin{tikzcd}[column sep=huge,row sep=large]
(X,\sigma) \arrow{r}{\text{Heisenberg}} 
  \arrow[bend right=20]{rr}[swap]{\text{CCR/Weyl}}
& \mathfrak{h}(X) \arrow{r}{\text{Hülle / Exponential}}
& U(\mathfrak{h}(X)),\ \mathrm{CCR}(X,\sigma)
\end{tikzcd}
\]

Diese Schicht beantwortet: \emph{Welche Observablen gibt es, und wie kommutieren sie?}

%--------------------------------------------------------------------
\emph{Schicht 3: Hilbertraum und Zustände -- Darstellungen}

\emph{Funktion dieser Schicht:}
Übergang von \emph{abstrakten} Observablen zu \emph{konkreten} Operatoren und Zuständen.

\emph{(1) Darstellung der CCR-Algebra.}

Eine (nichtdegenerierte) *-Darstellung
\[
\pi:\mathrm{CCR}(X,\sigma)\longrightarrow\mathcal{B}(\mathcal{H})
\]
ordnet jedem abstrakten Observable \(A\in\mathrm{CCR}(X,\sigma)\)
einen konkreten beschränkten Operator \(\pi(A)\) zu.

\begin{itemize}
  \item \(\mathcal{H}\) ist der \emph{Zustandsraum der Beschreibungen}.
  \item Vektoren \(\psi\in\mathcal{H}\) beschreiben Zustände (bis auf Phase);
        physikalische Zustände sind Strahlen \([\psi]\in\mathbb{P}(\mathcal{H})\).
  \item Erwartungswerte:
        \[
        \langle A\rangle_\psi := \langle\psi,\pi(A)\psi\rangle.
        \]
\end{itemize}

\emph{(2) Endliche QM als Spezialfall.}

Für \(X\cong\R^{2n}\):

\[
\mathcal{H}=L^2(\R^n),\quad \pi_{\mathrm{Sch}}(W(x,p))\psi(y)
= e^{\frac{i}{\hbar}(p\cdot y - \frac12 p\cdot x)}\psi(y-x).
\]

Die infinitesimalen Generatoren \(\hat X_j,\hat P_k\) der \(W(x,p)\) sind
\[
(\hat X_j\psi)(x)=x_j\psi(x),\quad
(\hat P_j\psi)(x)=-i\hbar\,\partial_{x_j}\psi(x).
\]

\medskip

\noindent
\textbf{Diagramm (kinematik + Zustände):}

\[
\begin{tikzcd}[column sep=huge,row sep=large]
(X,\sigma) \arrow{r}{\mathrm{CCR}}
& \mathcal{A}:=\mathrm{CCR}(X,\sigma) \arrow{r}{\pi}
& \pi(\mathcal{A})\subset\mathcal{B}(\mathcal{H})
\\
& & \mathcal{H} \arrow{u}[swap]{\psi\mapsto \langle\psi,\pi(\cdot)\psi\rangle}
\end{tikzcd}
\]

Diese Schicht beantwortet: \emph{Wie werden abstrakte Observablen als konkrete Operatoren realisiert, und was ist ein Zustand?}

%--------------------------------------------------------------------
\emph{Schicht 4: Symmetriegruppen $G$ und ihre Generatoren}

\emph{Funktion dieser Schicht:}
Kodierung von \emph{Symmetrien und Flüssen} (inklusive Zeitentwicklung) und ihrer Wirkung
auf Zustände und Observablen.

\emph{(1) Abstrakte Symmetrie auf der CCR-Ebene.}

Sei \(G\) eine Lie-Gruppe (Translationen, Rotationen, Zeitentwicklung, Gauge).
Wir verlangen eine \emph{Automorphismengruppe} der Observablenalgebra:

\[
\alpha:G\longrightarrow\mathrm{Aut}(\mathcal{A}),\qquad
g\mapsto\alpha_g,
\]
so dass
\[
\alpha_g:\mathcal{A}\to\mathcal{A}
\]
ein *-Automorphismus ist (Spektrum und Produktstruktur bleiben erhalten).

\emph{(2) Implementierung durch unitäre Darstellung.}

In einer Darstellung \(\pi\) soll \(\alpha\) durch unitäre Operatoren implementiert werden:

\[
\pi(\alpha_g(A)) = U(g)\,\pi(A)\,U(g)^{-1},\qquad U(g)\in\mathcal{U}(\mathcal{H}).
\]

Das ist genau die „Ad-Brücke“ aus deinem Lie-Gruppen-Text:
\[
dU(\mathrm{Ad}_g X) = \mathrm{Ad}_{U(g)}(dU(X)).
\]

\emph{(3) Generatoren als Observablen.}

Die zugehörige Lie-Algebra-Darstellung
\[
dU:\mathfrak g\to\mathfrak u(\mathcal{H})
\]
liefert die \emph{Generatoren}:
\[
\hat X := i\hbar\,dU(X),
\]
selbstadjungierte Operatoren, die physikalische Observablen sind (Impuls, Drehimpuls, Ladung, Energie).

\begin{itemize}
  \item Einparameteruntergruppe:
        \[
        g(t)=\exp(tX) \in G,\quad
        U(t):=U(g(t))=\exp\!\big(-\tfrac{i}{\hbar}t\hat X\big).
        \]
  \item Diese Flüsse realisieren \emph{Symmetrie- oder Zeitentwicklung} auf Zuständen:
        \(\psi(t)=U(t)\psi_0\), und auf Observablen: 
        \(A(t)=U(t)^{-1}AU(t)\).
\end{itemize}

\medskip

\noindent
\textbf{Diagramm (Symmetrieebene):}

\[
\begin{tikzcd}[column sep=huge,row sep=large]
G \arrow{r}{\alpha} \arrow{d}{U}
& \mathrm{Aut}(\mathcal{A})
\\
\mathcal{U}(\mathcal{H}) \arrow{ur}[swap]{\mathrm{Ad}}
\end{tikzcd}
\]

Infinitesimal:

\[
\begin{tikzcd}[column sep=huge,row sep=large]
\mathfrak g \arrow{r}{d\alpha} \arrow{d}{dU}
& \mathrm{Der}(\mathcal{A})
\\
\mathfrak u(\mathcal{H}) \arrow{ur}[swap]{\mathrm{ad}}
\end{tikzcd}
\]

wobei \(\mathrm{Der}(\mathcal{A})\) Derivationen der Algebra (z.\,B.\ Heisenberg-Gleichung)
und \(\mathrm{ad}_K(A)=[K,A]\) sind.

%--------------------------------------------------------------------
\emph{Schicht 5: Dynamik als spezieller Symmetriefluss}

\emph{Funktion dieser Schicht:}
Auswahl einer ausgezeichneten Einparametergruppe (Zeitentwicklung).

\begin{itemize}
  \item \textbf{Abstrakt (CCR-Ebene):}
        Eine Zeitentwicklungs-Automorphismengruppe
        \[
        \alpha_t:\mathcal{A}\to\mathcal{A},\quad t\in\R,
        \]
        mit \(\alpha_{t+s}=\alpha_t\circ\alpha_s\), \(\alpha_0=\mathrm{id}\).
        Infinitesimal:
        \[
        \frac{d}{dt}\alpha_t(A)\big|_{t=0} = \delta(A),
        \]
        wobei \(\delta\) eine Derivation ist (oft inner: \(\delta(A)=\frac{i}{\hbar}[H,A]\)).

  \item \textbf{Konkrete Realisierung:}
        In einer Darstellung \(\pi\) gibt es eine unitäre Einparametergruppe
        \[
        U(t)=e^{-\frac{i}{\hbar}t\hat H},\qquad \hat H \text{ selbstadjungiert},
        \]
        mit
        \[
        \pi(\alpha_t(A)) = U(t)\,\pi(A)\,U(t)^{-1}.
        \]

  \item \textbf{Zustandsentwicklung:}
        Schrödinger-Bild:
        \[
        i\hbar\frac{d}{dt}\psi(t) = \hat H \psi(t).
        \]
        Heisenberg-Bild:
        \[
        \frac{d}{dt}\hat A_H(t) = \frac{i}{\hbar}[\hat H,\hat A_H(t)].
        \]
\end{itemize}

\medskip

\noindent
\textbf{Gesamtbild (kommutatives Hauptdiagramm):}

\[
\begin{tikzcd}[column sep=huge,row sep=large]
(X,\sigma) \arrow{r}{\mathrm{CCR}} \arrow{d}{S_g\in\mathrm{Sp}(X,\sigma)}
& \mathcal{A} \arrow{r}{\pi} \arrow{d}{\alpha_g}
& \pi(\mathcal{A})\subset\mathcal{B}(\mathcal{H}) \arrow{d}{\mathrm{Ad}_{U(g)}}
\\
(X,\sigma) \arrow{r}{\mathrm{CCR}}
& \mathcal{A} \arrow{r}{\pi}
& \pi(\mathcal{A})
\end{tikzcd}
\]

\[
\begin{tikzcd}[column sep=huge,row sep=large]
\mathfrak g \arrow{r}{d\alpha} \arrow{d}{dU}
& \mathrm{Der}(\mathcal{A}) \arrow{d}{d\pi}
\\
\mathfrak u(\mathcal{H}) \arrow{r}{\mathrm{ad}}
& \mathrm{Der}(\pi(\mathcal{A}))
\end{tikzcd}
\]

Zeitentwicklung ist der Spezialfall, bei dem \(G=\R\), \(X\in\mathfrak g\simeq\R\)
und der Generator \(\hat H\) die Energie (Hamiltonoperator) ist.

%--------------------------------------------------------------------
\emph{Kurzinterpretation der Rollen}

\begin{REM}{TopDown-Funktionen}{Funktionsübersicht}
\begin{center}
\renewcommand{\arraystretch}{1.25}
\begin{tabular}{lll}
\textbf{Objekt} & \textbf{Ebene} & \textbf{Funktion} \\ \hline
\((X,\sigma)\) & klassisch & linearisierter Phasenraum; klassische Kinematik \\
\(\mathfrak{h}(X)\) & Lie-Algebra & infinitesimale CCR, nichtkommutative Kinematik \\
\(U(\mathfrak{h}(X))\) & assoz. Algebra & polynomiale Observablen, Heisenberg-Dynamik \\
\(\mathrm{CCR}(X,\sigma)\) & C$^\ast$-Algebra & abstrakte Observablenalgebra (beschränkt) \\
\(\pi:\mathrm{CCR}\to\mathcal{B}(\mathcal{H})\) & Darstellung & konkrete Operatoren, Zustandsbeschreibungen \\
\(\mathcal{H},\mathbb{P}(\mathcal{H})\) & Zustände & Beschreibungen / physikalische Zustände \\
\(G,\mathfrak g\) & Symmetriegruppe & Raum der Beschreibungstransformationen \\
\(U:G\to\mathcal{U}(\mathcal{H})\) & unitäre Darstellung & Wirkung von Symmetrien auf Zustände \\
\(\hat X=i\hbar dU(X)\) & Observablen & Generatoren von Flüssen (Impuls, Drehimpuls, Energie) \\
\(\alpha:G\to\mathrm{Aut}(\mathcal{A})\) & Automorphismen & Wirkung von Symmetrien auf Observablen \\
\(\hat H\) & Hamiltonian & Generator der Zeitentwicklung
\end{tabular}
\end{center}

Konzeptionell:

\[
\boxed{
\begin{aligned}
&\text{CCR-/Heisenberg-Struktur:}\;\;\text{„Was sind die Observablen, wie kommutieren sie?“}\\
&\text{Hilbertraum + Darstellung:}\;\;\text{„Wie sehen sie als Operatoren aus, was ist ein Zustand?“}\\
&\text{Lie-Gruppen + Flüsse:}\;\;\text{„Wie transformieren Zustände und Observablen (Symmetrien, Zeit)?“}
\end{aligned}
}
\]
\end{REM}


\end{unit}





\pagebreak

\section{CPT Transformations}

\subsection{Systematische Übersicht der diskreten Symmetrien $\mathcal P$, $\mathcal C$, $\mathcal T$}

\begin{MOT}{QFT.CPT.S1}{Einordnung und Ziel}
Die drei diskreten Transformationen \emph{Parität} ($\mathcal P$), \emph{Ladungskonjugation} ($\mathcal C$) und \emph{Zeitumkehr} ($\mathcal T$) bilden die fundamentalen diskreten Symmetrien der relativistischen Quantenfeldtheorie.  
Sie wirken auf unterschiedliche Aspekte eines Feldes:
\begin{itemize}
  \item $\mathcal P$: Raumspiegelung — betrifft die räumlichen Koordinaten.
  \item $\mathcal C$: Konjugation der inneren (Ladungs-)Struktur.
  \item $\mathcal T$: Umkehr der Zeitrichtung — betrifft Dynamik und Kausalordnung.
\end{itemize}
Obwohl sie auf unterschiedlichen Ebenen wirken, sind alle drei Transformationen durch die Forderung der \emph{Kovarianz der Diracgleichung} eindeutig bestimmt.
\end{MOT}

\begin{CONC}{QFT.CPT.S2}{Vergleichende Übersicht}

\renewcommand{\arraystretch}{1.4}
\begin{center}
\begin{tabular}{|c|c|c|c|}
\hline
\textbf{Transformation} & \textbf{Definition im Raumzeitpunkt} & \textbf{Spinortransformation} $S(\mathcal X)$ & \textbf{Feldwirkung} $\psi^{\mathcal X}(x)$ \\ \hline
\textbf{Parität} $\mathcal P$ &
$x'=(t,-\mathbf x)$ &
$S(\mathcal P)=\gamma^0$ &
$\psi^{\mathcal P}(t,\mathbf x)=\gamma^0\,\psi(t,-\mathbf x)$ \\ \hline
\textbf{Ladungskonjugation} $\mathcal C$ &
$x'=x$ (keine Raumzeitänderung) &
$S(\mathcal C)=i\gamma^2\gamma^0$ &
$\psi^{\mathcal C}(x)=S(\mathcal C)\,\bar\psi^{\,T}(x)=i\gamma^2\gamma^0\bar\psi^{\,T}(x)$ \\ \hline
\textbf{Zeitumkehr} $\mathcal T$ &
$x'=(-t,\mathbf x)$ &
$S(\mathcal T)=\gamma_5\gamma^0$ &
$\psi^{\mathcal T}(t,\mathbf x)=\gamma^1\gamma^3\,\psi^*(-t,\mathbf x)$ \\ \hline
\end{tabular}
\end{center}

\end{CONC}

\begin{CONC}{QFT.CPT.S3}{Transformationsgesetze der $\gamma$-Matrizen}
Die Bedingung der \emph{Kovarianz der Diracgleichung}
\[
S^{-1}(\mathcal X)\,\gamma^\mu\,S(\mathcal X)=\Lambda^\mu{}_\nu(\mathcal X)\,\gamma^\nu
\]
führt für jede diskrete Transformation auf charakteristische Relationen:

\begin{center}
\renewcommand{\arraystretch}{1.4}
\begin{tabular}{|c|c|c|c|}
\hline
$\mathcal X$ & $\Lambda(\mathcal X)$ & $S^{-1}\gamma^0S$ & $S^{-1}\gamma^iS$ \\ \hline
$\mathcal P$ & $\mathrm{diag}(1,-1,-1,-1)$ & $\gamma^0$ & $-\gamma^i$ \\ \hline
$\mathcal C$ & $\mathrm{diag}(1,1,1,1)$ (intern) & $-(\gamma^0)^T$ & $-(\gamma^i)^T$ \\ \hline
$\mathcal T$ & $\mathrm{diag}(-1,1,1,1)$ & $-\gamma^0$ & $+\gamma^i$ \\ \hline
\end{tabular}
\end{center}

Damit sind die Matrizen $S(\mathcal X)$ durch die jeweiligen Forderungen eindeutig (bis auf Phasen) bestimmt:
\[
S(\mathcal P)=\gamma^0,\quad
S(\mathcal C)=i\gamma^2\gamma^0,\quad
S(\mathcal T)=\gamma_5\gamma^0.
\]
\end{CONC}

\begin{EXA}{QFT.CPT.S4}{Wirkung auf Basis-Spinoren $u_s(p),v_s(p)$}
\renewcommand{\arraystretch}{1.5}
\begin{center}
\begin{tabular}{|c|c|c|}
\hline
Transformation & $u_s(p)$ & $v_s(p)$ \\ \hline
$\mathcal P$ &
$\gamma^0u_s(p')=+u_s(p)$ &
$\gamma^0v_s(p')=-v_s(p)$ \\ \hline
$\mathcal C$ &
$S(\mathcal C)\bar u_s^{\,T}(p)=(-1)^{s+1}v_s(p)$ &
$S(\mathcal C)\bar v_s^{\,T}(p)=(-1)^{s+1}u_s(p)$ \\ \hline
$\mathcal T$ &
$\gamma^1\gamma^3u_s^*(p)=(-1)^s u_r(E,-\mathbf p)$ &
$\gamma^1\gamma^3v_s^*(p)=(-1)^{s+1} v_r(E,-\mathbf p)$ \\ \hline
\end{tabular}
\end{center}

\noindent
Hierbei gilt $p'=(E,-\mathbf p)$ und $r\neq s$.  
Die Phasenwahl in $\mathcal P,\mathcal C,\mathcal T$ ist konventionsabhängig, hat jedoch keine physikalischen Konsequenzen.
\end{EXA}

\begin{REM}{QFT.CPT.S5}{Mathematische Struktur}
\begin{itemize}
  \item $\mathcal P$ und $\mathcal T$ sind \emph{uneigentliche Lorentztransformationen} und wirken auf die Raumzeitkoordinaten.  
  Im Spinorraum entsprechen sie Elementen der $\Pin(1,3)$-Gruppe.
  \item $\mathcal C$ ist eine \emph{interne, unitäre Symmetrie}, die den Ladungsoperator $Q$ in $-Q$ überführt:
  \[
  \mathcal C Q \mathcal C^{-1} = -Q.
  \]
  \item $\mathcal T$ ist \emph{antiunitär}, d.\,h.\ sie wirkt antilinear:
  \[
  \mathcal T\,i\,\mathcal T^{-1}=-i.
  \]
\end{itemize}
\end{REM}

\begin{CONC}{QFT.CPT.S6}{Physikalische Interpretation (Parallelisierung)}
\begin{itemize}
  \item \textbf{Parität} $\mathcal P$: Raumspiegelung — kehrt Impulsrichtung $\mathbf p\!\to\!-\mathbf p$ um, lässt Energie und Spinbetrag unverändert.  
    Intrinsische Parität trennt Teilchen ($+$) und Antiteilchen ($-$).
  \item \textbf{Ladungskonjugation} $\mathcal C$: Austausch von Teilchen und Antiteilchen — vertauscht Operatoren $a_s(\mathbf p)\leftrightarrow b_s(\mathbf p)$, koppelt an entgegengesetzte elektromagnetische Felder.
  \item \textbf{Zeitumkehr} $\mathcal T$: Inversion der Zeitrichtung $t\!\to\!-t$ — Impuls $\mathbf p\!\to\!-\mathbf p$, Spinrichtung kehrt sich um ($s\leftrightarrow r$), Operation ist antiunitär.
\end{itemize}
\end{CONC}

\begin{THEO}{QFT.CPT.S7}{Gemeinsame Invarianz — CPT-Theorem (Ausblick)}
Unter sehr allgemeinen Bedingungen (Lokalität, Spektralbedingung, Lorentzkovarianz, eindeutiges Vakuum) gilt:
\[
\boxed{(i\gamma^\mu\partial_\mu - m)\psi(x)=0
\quad\Longrightarrow\quad
(i\gamma^\mu\partial_\mu - m)\psi^{\mathcal{CPT}}(-x)=0.}
\]
Die kombinierte Transformation
\[
\psi^{\mathcal{CPT}}(x)
= S(\mathcal T)S(\mathcal P)S(\mathcal C)\,\psi^*(-x)
\]
ist eine exakte Symmetrie aller lokal-relativistischen Quantenfeldtheorien.  
Die CPT-Invarianz garantiert u.a. Massen- und Lebensdauer-Gleichheit von Teilchen und Antiteilchen und verknüpft fundamentale Symmetrien von Raum, Zeit und Ladung.
\end{THEO}




\section{Archiv}

\begin{unit}{}
\emph{Abstrakte Heisenberg- und Weyl-Struktur}

\emph{1. Abstrakter symplektischer Phasenraum}

Sei \(X\) ein reeller lokalkonvexer Vektorraum (typisch: endlichdimensional
\(X\simeq\R^{2n}\) oder ein Raum von Testfunktionen), versehen mit einer
bilinearen, antisymmetrischen Form
\[
\sigma : X\times X \longrightarrow \R,
\qquad \sigma(f,g) = -\sigma(g,f),
\]
die nichtausgeartet ist. Das Paar \((X,\sigma)\) heißt \emph{(abstrakter) symplektischer Raum}.

\begin{EXA}{HEIS-ABS.1}{Beispiele}
\begin{enumerate}
\item \emph{Endlichdimensional (QM):}
      \(X = V_x\oplus V_p\simeq\R^{2n}\) mit
      \(\sigma((x,p),(x',p')) = p\cdot x' - p'\cdot x\).
\item \emph{Feldtheorie (freie bosonische Felder):}
      \(X\) ist ein Raum reeller Testfunktionen,
      z.\,B. \(X=\mathcal{S}(\R^d)\),
      und \(\sigma\) ist durch ein zweipoliges Kernel gegeben
      (z.\,B. das Pauli--Jordan-Funktional).
\end{enumerate}
\end{EXA}

\emph{2. Abstrakte Heisenberg-Lie-Algebra}

\begin{DEF}{HEIS-ABS.2}{Abstrakte Heisenberg-Algebra}
Zu einem symplektischen Raum \((X,\sigma)\) definieren wir die
\emph{abstrakte Heisenberg-Lie-Algebra}
\[
\mathfrak{h}(X) := X \oplus \R Z
\]
mit Lie-Klammer
\[
[(f,\alpha Z),(g,\beta Z)]
:= \bigl(0,\ \sigma(f,g)\,Z\bigr),
\]
d.\,h.
\[
[f,g] = \sigma(f,g)\,Z,\qquad [Z,\cdot]=0.
\]
\end{DEF}

\begin{REM}{HEIS-ABS.3}{Interpretation}
\begin{itemize}
\item Für \(X=\R^{2n}\) mit kanonischer symplektischer Form
      reproduziert \(\mathfrak{h}(X)\) die übliche Heisenberg-Algebra
      mit endlich vielen \(X_j,P_k\).
\item Für \(X\) als Raum von Testfunktionen entsteht eine
      \emph{unendlichdimensionale} Heisenberg-Algebra, deren Generatoren
      man physikalisch als „abstrakte Feldmoden“ interpretieren kann:
      \(f\mapsto \Phi(f)\) mit
      \([\,\Phi(f),\Phi(g)\,]= i\,\sigma(f,g)\,\id\).
\end{itemize}
\end{REM}

Die zugehörige \emph{universelle Hüllalgebra} \(U(\mathfrak{h}(X))\)
ist dann die assoziative (algebraische) Umgebung aller polynomiellen
Kombinationen solcher „abstrakter Observablen“.

\emph{3. Abstrakte Weyl-Algebra (CCR-Algebra)}

\begin{DEF}{Weyl-ABS.1}{Abstrakte Weyl- (CCR-) Algebra}
Sei \((X,\sigma)\) wie oben.
Die \emph{Weyl-Algebra} \(\mathcal{W}(X,\sigma)\) ist die
assoziative *-Algebra, erzeugt von formalen Symbolen \(W(f)\) (\(f\in X\)),
unter den \emph{Weyl-Relationen}
\begin{align*}
W(f)W(g) &= e^{-\frac{i}{2}\sigma(f,g)}\,W(f+g),\\
W(f)^\ast &= W(-f),\\
W(0) &= \mathbf{1}.
\end{align*}
\end{DEF}

Man kann daraus durch Vervollständigung eine C\(^\ast\)-Algebra
\(\mathrm{CCR}(X,\sigma)\) konstruieren
(\emph{Weyl--C\(^\ast\)-Algebra} oder \emph{CCR-Algebra}),
die als abstrakte Observablenalgebra der bosonischen Theorie dient.

\begin{EXA}{Weyl-ABS.2}{QM und QFT als Spezialfälle}
\begin{enumerate}
\item \emph{QM (endlichdimensional):}
      \(X=\R^{2n}\) mit kanonischer \(\sigma\).
      Die Weyl-Algebra \(\mathcal{W}(X,\sigma)\) ist
      (bis auf Isomorphie) die bekannte Weyl-Algebra der Operatoren
      \(e^{\frac{i}{\hbar}(p\cdot X - x\cdot P)}\).
      Nach dem Satz von Stone--von Neumann sind alle irreduziblen
      Darstellungen (unter milden Bedingungen) unitar äquivalent.
\item \emph{Freie bosonische QFT:}
      \(X\) ist ein (realer) Raum von Lösungen der klassischen Feldgleichung
      modulo Nebenbedingungen, \(\sigma\) ist die entsprechende
      symplektische Form (z.\,B. durch das Pauli--Jordan-Funktional).
      Die Weyl-C\(^\ast\)-Algebra \(\mathrm{CCR}(X,\sigma)\) ist die
      abstrakte Observablenalgebra des freien Feldes.
      Es gibt nun \emph{viele} unitär nichtäquivalente Darstellungen,
      die jeweils unterschiedliche physikalische Vakuumzustände,
      Temperaturzustände oder Hintergrundgeometrien repräsentieren.
\end{enumerate}
\end{EXA}

\emph{4. Dynamik als Automorphismengruppe}

Auf dieser Ebene kann man Dynamik abstrakt fassen als
eindimensionale Gruppen von *-Automorphismen:
\begin{DEF}{DYN-ABS.1}{Abstrakte Dynamik}
Eine (Heisenberg-)Dynamik auf \(\mathcal{W}(X,\sigma)\) ist
eine stark-stetige Gruppe von *-Automorphismen
\[
\alpha : \R \to \mathrm{Aut}(\mathcal{W}(X,\sigma)),\qquad
t\mapsto \alpha_t,
\]
die (formal) durch einen „Hamiltonoperator“ \(H\) implementiert wird:
\[
\frac{d}{dt}\alpha_t(A)\big|_{t=0}
= \delta(A) = i[H,A] \quad \text{(wo sinnvoll definiert)}.
\]
\end{DEF}

In einer konkreten Darstellung \(\pi:\mathcal{W}(X,\sigma)\to\mathcal{B}(\mathcal{H})\)
wird diese Dynamik durch eine unitäre Gruppe \(U(t)\) implementiert:
\[
\pi(\alpha_t(A)) = U(t)\,\pi(A)\,U(t)^\ast.
\]

\emph{5. Kritische Bewertung der Allgemeinheit}

\begin{REM}{GEN-ABS.3}{Reichweite und Grenzen}
\begin{itemize}
\item Die abstrakte Konstruktion \((X,\sigma)\mapsto \mathfrak{h}(X)\mapsto \mathcal{W}(X,\sigma)\)
      erfasst \emph{alle bosonischen Theorien mit kanonischen Kommutationsrelationen}:
      \begin{itemize}
      \item gewöhnliche QM (endlichdimensionales \(X\)),
      \item freie bosonische QFT (unendlichdimensionales \(X\)).
      \end{itemize}
\item Fermionische Theorien (Dirac-Feld) erfordern statt \((X,\sigma)\)
      einen reellen Hilbertraum \((Y,\langle\cdot,\cdot\rangle)\)
      und die CAR-Algebra \(\mathrm{CAR}(Y)\); formal ist dies
      die „Clifford-Version“ der obigen Konstruktion.
\item Allgemeine quantenmechanische Systeme können Observablenalgebren
      besitzen, die \emph{nicht} durch CCR/CAR beschreibbar sind
      (z.\,B. reine Spin-Systeme mit \(\mathfrak{su}(2)\)-Algebra,
      quantisierte Gauge-Theorien mit Nebenbedingungen etc.).
      In voller Allgemeinheit arbeitet man daher mit einer
      abstrakten C\(^\ast\)- oder von-Neumann-Algebra \(\mathcal{A}\)
      und deren Darstellungen.
\end{itemize}

\emph{Fazit.}
Die abstrakte Heisenberg-/Weyl-Formalisierung über \((X,\sigma)\)
liefert einen einheitlichen Rahmen, in dem
\begin{itemize}
\item endliche QM und freie bosonische QFT
      als Spezialfälle desselben kinematischen Prinzips erscheinen,
\item Dynamik als Automorphismengruppe der Observablenalgebra formuliert wird,
\item und Zustände erst in einem zweiten Schritt als positive Funktionale
      auf dieser Algebra eingeführt werden.
\end{itemize}
Für eine vollständig allgemeine Charakterisierung \emph{aller} quantenmechanischen
Theorien reicht sie jedoch nicht aus; dort benötigt man
den abstrakten Rahmen der Operatoralgebren, in dem die CCR/CAR/Weyl-Algebren
wichtige, aber nicht die einzigen Beispiele sind.
\end{REM}
\end{unit}

\begin{unit}
\emph{Abstrakter Observablen-Formalismus via C$^\ast$- und von-Neumann-Algebren}

In diesem Abschnitt heben wir den zuvor beschriebenen Heisenberg-/Weyl-Formalismus
auf eine noch allgemeinere Ebene:
Statt CCR/CAR-spezifischer Algebren betrachten wir \emph{abstrakte Operatoralgebren}
(C$^\ast$- und von-Neumann-Algebren) als Träger der Quantenmechanik.
Die Heisenberg- und Weyl-Algebren erscheinen dann als wichtige Spezialfälle.

%---------------------------------------------------------
\emph{1. C$^\ast$-Algebren als abstrakte Observablenalgebren}

\begin{DEF}{CSTAR-1}{C$^\ast$-Algebra}
Eine \emph{C$^\ast$-Algebra} ist eine komplexe Banach-Algebra \((\mathcal{A},\|\cdot\|)\)
mit involutiver Antilinearität \(^{\ast}:\mathcal{A}\to\mathcal{A}\),
die die folgenden Eigenschaften erfüllt:
\begin{enumerate}
  \item \((ab)^\ast = b^\ast a^\ast,\quad (\lambda a + \mu b)^\ast = \overline{\lambda}a^\ast + \overline{\mu}b^\ast\).
  \item \(\|a^\ast a\| = \|a\|^2\) für alle \(a\in\mathcal{A}\).
  \item \(\mathcal{A}\) ist vollständig bezüglich der Norm \(\|\cdot\|\).
\end{enumerate}
\end{DEF}

\begin{REM}{CSTAR-2}{Interpretation als Observablenalgebra}
\begin{itemize}
  \item Die selbstadjungierten Elemente \(a=a^\ast\in\mathcal{A}\) werden als \emph{Observablen} interpretiert.
  \item Die Algebra selbst beschreibt die \emph{vollständige Menge} aller beobachtbaren Größen
        und ihrer algebraischen Relationen (Produkt, Spektrum, funktionale Kalküle).
  \item Eine \emph{konkrete Darstellung} \(\pi:\mathcal{A}\to\mathcal{B}(\mathcal{H})\) auf einem Hilbertraum \(\mathcal{H}\)
        ist nur eine (mögliche) Realisierung; die abstrakte Algebra \(\mathcal{A}\) ist primär.
\end{itemize}
\end{REM}

\begin{EXA}{CSTAR-3}{Wichtige Beispiele}
\begin{enumerate}
  \item \emph{Quantenmechanik mit endlich vielen Freiheitsgraden:}
        \(\mathcal{A} = \mathcal{B}(\mathcal{H})\), die Algebra der beschränkten Operatoren
        auf einem separablen Hilbertraum \(\mathcal{H}\) (z.\,B. \(\mathcal{H}=L^2(\R^n)\)).
  \item \emph{CCR- und CAR-Algebren:}
        Für einen symplektischen Raum \((X,\sigma)\) ist \(\mathrm{CCR}(X,\sigma)\) 
        die Weyl-C$^\ast$-Algebra; für einen reellen Hilbertraum \(Y\) ist \(\mathrm{CAR}(Y)\)
        die CAR-C$^\ast$-Algebra. Beide sind spezielle C$^\ast$-Algebren.
  \item \emph{Kommutative C$^\ast$-Algebren:}
        Jede kommutative C$^\ast$-Algebra ist isomorph zu \(C_0(M)\) für einen
        lokalkompakten Hausdorffraum \(M\) (Gelfand-Naimark-Theorem).
        Dies liefert die Brücke zur klassischen Mechanik: Phasenraum \(M\) $\Leftrightarrow$ kommutative Observablenalgebra \(C_0(M)\).
\end{enumerate}
\end{EXA}

%---------------------------------------------------------
\emph{2. Zustände und GNS-Konstruktion}

\begin{DEF}{STATE-1}{Zustand auf einer C$^\ast$-Algebra}
Ein \emph{Zustand} auf einer C$^\ast$-Algebra \(\mathcal{A}\) ist ein lineares Funktional
\(\omega:\mathcal{A}\to\C\), das folgende Eigenschaften erfüllt:
\begin{enumerate}
  \item \emph{Positivität:} \(\omega(a^\ast a)\ge0\) für alle \(a\in\mathcal{A}\).
  \item \emph{Normierung:} \(\omega(\mathbf{1})=1\) (falls \(\mathcal{A}\) eine Eins besitzt).
\end{enumerate}
\end{DEF}

\begin{REM}{STATE-2}{Physikalische Bedeutung}
\begin{itemize}
  \item Der Erwartungswert einer Observablen \(A=A^\ast\in\mathcal{A}\) im Zustand \(\omega\)
        ist \(\omega(A)\).
  \item Zustände sind also die abstrakte Form dessen, was in konkreten Darstellungen
        durch Vektoren \(\psi\in\mathcal{H}\) oder Dichtematrizen \(\rho\) beschrieben wird.
\end{itemize}
\end{REM}

\begin{THEO}{GNS-1}{GNS-Konstruktion (Gelfand--Naimark--Segal)}
Zu jedem Zustand \(\omega\) auf einer C$^\ast$-Algebra \(\mathcal{A}\) existiert ein Tripel
\((\pi_\omega,\mathcal{H}_\omega,\Omega_\omega)\), bestehend aus
\begin{enumerate}
  \item einer Hilbertraumdarstellung \(\pi_\omega:\mathcal{A}\to\mathcal{B}(\mathcal{H}_\omega)\),
  \item einem zyklischen Vektor \(\Omega_\omega\in\mathcal{H}_\omega\)
        (d.\,h. \(\pi_\omega(\mathcal{A})\Omega_\omega\) ist dicht in \(\mathcal{H}_\omega\)),
\end{enumerate}
so dass
\[
\omega(A) = \langle \Omega_\omega,\,\pi_\omega(A)\Omega_\omega\rangle_{\mathcal{H}_\omega}
\quad\text{für alle }A\in\mathcal{A}.
\]
Dieses Tripel ist bis auf unitären Isomorphismus eindeutig.
\end{THEO}

\begin{REM}{GNS-2}{Interpretation}
\begin{itemize}
  \item Die GNS-Konstruktion zeigt: \emph{Jeder abstrakte Zustand definiert seinen eigenen Hilbertraum}
        und eine Darstellung der Observablenalgebra.
  \item Dies ist in der QFT essentiell: unterschiedliche Vakuumzustände (Minkowski, Rindler, Thermalzustände)
        führen i.\,Allg. zu nicht-unitär äquivalenten Darstellungen.
  \item Die „Welt“ eines Beobachters ist also \((\pi_\omega(\mathcal{A}),\mathcal{H}_\omega,\Omega_\omega)\).
\end{itemize}
\end{REM}

%---------------------------------------------------------
\emph{3. Dynamik und Symmetrien als Automorphismen}

\begin{DEF}{DYN-ABS}{Automorphismen und Dynamik}
Sei \(\mathcal{A}\) eine C$^\ast$-Algebra.
\begin{enumerate}
  \item Ein *-Automorphismus \(\alpha:\mathcal{A}\to\mathcal{A}\) ist ein bijektiver *-Homomorphismus.
  \item Eine \emph{(Heisenberg-)Dynamik} ist eine stark-stetige Einparametergruppe
        \(\alpha:\R\to\mathrm{Aut}(\mathcal{A})\), \(t\mapsto\alpha_t\):
        \[
        \alpha_0=\id,\quad \alpha_{t+s}=\alpha_t\circ\alpha_s,\quad
        t\mapsto\alpha_t(A)\text{ stetig für alle }A\in\mathcal{A}.
        \]
\end{enumerate}
\end{DEF}

\begin{REM}{DYN-IMP}{Physikalische Bedeutung}
\begin{itemize}
  \item \(\alpha_t\) beschreibt die Zeitentwicklung der Observablen im Heisenberg-Bild: 
        \(\alpha_t(A)\) ist die Observable „\(A\) zur Zeit \(t\)“.
  \item In einer GNS-Darstellung \((\pi_\omega,\mathcal{H}_\omega,\Omega_\omega)\)
        wird \(\alpha_t\) oft (unter geeigneten Bedingungen) durch eine unitäre Gruppe \(U_\omega(t)\)
        implementiert:
        \[
        \pi_\omega(\alpha_t(A)) = U_\omega(t)\,\pi_\omega(A)\,U_\omega(t)^\ast.
        \]
  \item Analog werden Symmetriegruppen (z.\,B. Poincaré-Gruppe) durch
        Automorphismengruppen \(\beta_g\in\mathrm{Aut}(\mathcal{A})\) repräsentiert,
        die in Darstellungen durch unitäre (oder projektiv unitäre) Gruppen realisiert werden.
\end{itemize}
\end{REM}

%---------------------------------------------------------
\emph{4. von-Neumann-Algebren und lokale Strukturen}

\begin{DEF}{VN-1}{von-Neumann-Algebra}
Eine \emph{von-Neumann-Algebra} ist eine *-Unteralgebra
\(\mathcal{M}\subset \mathcal{B}(\mathcal{H})\), die
\begin{enumerate}
  \item schwach operator-topologisch abgeschlossen ist,
  \item die Eins \(\mathbf{1}\) enthält.
\end{enumerate}
Äquivalent: \(\mathcal{M} = \mathcal{M}''\), d.\,h. \(\mathcal{M}\) ist gleich ihrem Bicommutant.
\end{DEF}

\begin{REM}{VN-2}{Rolle in QFT}
\begin{itemize}
  \item In der algebraischen QFT (Haag–Kastler-Formalismus) ordnet man jeder
        (geeigneten) Raumzeitregion \(\mathcal{O}\) eine von-Neumann-Algebra
        \(\mathcal{A}(\mathcal{O})\subset\mathcal{B}(\mathcal{H})\) zu, die die
        \emph{lokalen Observablen in \(\mathcal{O}\)} beschreibt.
  \item Die Gesamtheit \(\{\mathcal{A}(\mathcal{O})\}_{\mathcal{O}}\) bildet ein
        \emph{lokales Netz von Operatoralgebren}, das Kovarianz- und Lokalitätsaxiome erfüllt.
  \item Die Struktur solcher Netze (Typ von von-Neumann-Algebren, Superselektionssektoren, etc.)
        trägt tiefe Informationen über die Feldtheorie.
\end{itemize}
\end{REM}

%---------------------------------------------------------
\emph{5. Einordnung der CCR/CAR/Weyl-Algebren in diesen Rahmen}

\begin{EXA}{ABSTR-CCR-1}{CCR/CAR/Weyl als Spezialfälle}
\begin{itemize}
  \item Für einen symplektischen Raum \((X,\sigma)\) ist
        \(\mathrm{CCR}(X,\sigma)\) eine C$^\ast$-Algebra,
        die die Weyl-Relationen kodiert.
        Sie ist ein Spezialfall des allgemeinen C$^\ast$-Formalismus.
  \item Für einen reellen Hilbertraum \(Y\) ist
        \(\mathrm{CAR}(Y)\) die C$^\ast$-Algebra der Canonical Anti-commutation Relations,
        zentral für fermionische Systeme.
  \item Die früher eingeführten Heisenberg- und Weyl-Algebren
        können als \emph{dichte *-Unteralgebren} dieser C$^\ast$-Algebren
        aufgefasst werden (auf algebraischer Ebene).
\end{itemize}
\end{EXA}

\begin{REM}{ABSTR-CCR-2}{Abstrakter Oberbegriff}
\begin{itemize}
  \item Der C$^\ast$- bzw. von-Neumann-Formalismus liefert einen \emph{Oberbegriff}
        für quantenmechanische Theorien: Eine Theorie ist charakterisiert durch
        \[
        (\mathcal{A},\mathcal{S},\alpha_t,G),
        \]
        wobei \(\mathcal{A}\) eine C$^\ast$- oder W$^\ast$-Algebra von Observablen ist,
        \(\mathcal{S}\) eine Menge zulässiger Zustände,
        \(\alpha_t\) eine Dynamik (Automorphismengruppe)
        und \(G\) eine Symmetriegruppe mit Automorphismen \(\beta_g\).
  \item CCR-/CAR-/Weyl-Algebren sind dann konkrete Beispiele solcher \(\mathcal{A}\),
        insbesondere für feldtheoretische Modelle mit kanonischen Relationen.
  \item Spin-Systeme, Lattice-Modelle, topologische Quantensysteme etc.\ lassen sich
        ebenfalls in diesen Rahmen integrieren, ohne CCR/CAR-Struktur vorauszusetzen.
\end{itemize}
\end{REM}

%---------------------------------------------------------
\emph{6. Zusammenfassung: Allgemeinheit und Stringenz}

\begin{REM}{ABSTR-SUM}{Reichweite des C$^\ast$/vN-Standpunkts}
\begin{itemize}
  \item \textbf{Formale Stringenz:}
        Der C$^\ast$-algebraische Zugang ist vollständig funktionalanalytisch fundiert;
        GNS sichert die Verbindung zur üblichen Hilbertraum-Formulierung.
  \item \textbf{Allgemeinheit:}
        Er umfasst sowohl
        \begin{itemize}
          \item gewöhnliche Quantenmechanik (endlich viele Freiheitsgrade),
          \item freie und (in geeigneten Konstruktionen) interagierende QFT,
          \item nicht-kanonische Systeme (Spin, Lattice-Modelle, topologische Theorien).
        \end{itemize}
  \item \textbf{Einordnung der Heisenberg-/Weyl-Struktur:}
        Die zuvor eingeführten Heisenberg-, Hüll- und Weyl-Algebren
        sind in diesem Rahmen als besonders strukturierte Beispiele von
        Observablenalgebren zu verstehen, die zusätzlich eine klare
        symplektische Herkunft (CCR) haben.
        Sie liefern einen sehr transparenten kinematischen Kern,
        aber nicht die einzig mögliche Algebra-Architektur.
\end{itemize}

\emph{Fazit.}
Während die abstrakte Heisenberg-/Weyl-Formalisierung über \((X,\sigma)\)
bosonische Systeme mit kanonischen Relationen sehr elegant erfasst,
bietet der C$^\ast$/von-Neumann-Formalismus eine noch allgemeinere,
operatoralgebraische Bühne, in der „Quantenmechanik“ im weitesten Sinn
als Struktur aus \emph{Observablenalgebra, Zuständen, Dynamik und Symmetrien}
definiert werden kann.
\end{REM}
\end{unit}

\pagebreak
\printbibliography
\end{document}